\documentclass[11pt,letterpaper]{article}
\usepackage[margin=1in]{geometry}
\usepackage[T1]{fontenc}
\usepackage{lmodern}
\usepackage{amsmath,amssymb,amsthm,mathtools}
\usepackage{booktabs,array,tabularx}
\usepackage{microtype}
\usepackage[numbers,sort&compress]{natbib}
\usepackage{xcolor}
\usepackage[colorlinks=true,linkcolor=blue!45!black,citecolor=blue!45!black,urlcolor=blue!45!black]{hyperref}
\hypersetup{pdftitle={Positional Information for Budget-Adaptive Markov Prediction},pdfauthor={},pdfsubject={Local research draft}}
\newcommand{\TV}{\operatorname{TV}}
\newcommand{\E}{\mathbb E}
\newcommand{\ind}{\mathbf 1}
\newcommand{\Prb}{\mathbb P}
\newcommand{\eps}{\varepsilon}
\newcommand{\cT}{\mathcal T}
\newcommand{\cK}{\mathcal K}
\newcommand{\R}{\mathbb R}
\DeclareMathOperator{\diag}{diag}
\DeclareMathOperator{\tr}{tr}
\DeclareMathOperator{\Var}{Var}
\DeclareMathOperator{\Cov}{Cov}
\DeclareMathOperator{\KL}{KL}
\newtheorem{theorem}{Theorem}[section]
\newtheorem{proposition}[theorem]{Proposition}
\newtheorem{lemma}[theorem]{Lemma}
\newtheorem{corollary}[theorem]{Corollary}
\theoremstyle{definition}
\newtheorem{definition}[theorem]{Definition}
\theoremstyle{remark}
\newtheorem{remark}[theorem]{Remark}
\numberwithin{equation}{section}
\setlength{\emergencystretch}{2em}
\title{Positional Information for Budget-Adaptive Markov Prediction}
\author{}
\date{Local finite-horizon research draft, 2 November 2026}

\begin{document}
\maketitle
\begin{abstract}
Consider prediction from one binary order-two Markov trajectory after an unknown number of symbol edits, with the terminal pair protected. The loss is expected total variation at that actual pair. At a supplied all-state two-step minorization certificate $1/M$, the known-budget risk scale is $\max\{M^{-1/2},k/M\}$. We ask which restricted observations preserve this scale simultaneously for every budget $k$. For every finite horizon $M\ge64$ and integer $2\le r\le M+2$, the optimal budget-blind competitive factor from exact triple counts and the first and last $r$ symbols lies between $2^{-30}\ell$ and $2^{20}\ell$, where $\ell=1+\log((M+2)/r)$. The upper uses only counts and the suffix; the lower permits both windows and compares the same map across budgets, rather than imposing a logarithmic penalty at every budget. Consequently a prescribed width sequence admits a finite uniform factor exactly when $\liminf_M r_M/M>0$. A separate fixed-fraction argument retains stronger explicit lower constants. The positive construction accepts an empirical row estimate only when enough visits occurred before the observed suffix. The lower combines an occupation mixture with a count-preserving interior rearrangement and a pointwise legal edit kernel. Under a different access model, a whole-word scan can compress the needed gate into one exact supplemental bit. These are statistical observation results, not learning or architectural separation claims.
\end{abstract}

\section{The observation question}
A transition histogram records which local patterns occurred but not when they occurred. For a clean homogeneous Markov trajectory, this is consistent with the usual likelihood factorization into transition counts and the initial state. With position-sensitive corruption, however, the observation channel changes. A predictor may need to distinguish an empirical row supported by older visits from one populated almost entirely near the prediction time. The question here is how much positional information is needed when the number of edits is not supplied.

We study one concrete experiment: a binary order-two source, $M$ observed departures, and an unmodified terminal pair. The source obeys an all-state two-step minorization with supplied coefficient $1/M$. An offline kernel may change at most $k$ earlier symbols on every realization. The predictor sees exact triple counts and prescribed boundary windows, but not $k$. Its risk is expected unsquared total variation to the next-symbol distribution at the actual terminal pair. In this experiment, giving $k$ to the predictor makes counts and the protected pair sufficient for the optimal \emph{rate}; no claim of statistical sufficiency is needed.

Our main result characterizes the competitive factor simultaneously for every finite horizon and boundary width: its order is $1+\log((M+2)/r)$ with universal constants. One map at each fixed horizon must serve every budget. The positive lower limiting fraction criterion and logarithmic fixed-fraction price follow at the order level; a separate fixed-fraction proof supplies a stronger explicit lower constant. A separate corollary shows why raw access and computed messages must not be identified: a whole-word scan can produce a single exact gate bit that suffices beyond counts. The scan, exact counters, and protected query are part of that corollary's information model.

The proof mechanisms are complementary. For the upper bound, a gate retains evidence of visits before a suffix, removing the shortest-return terms from a harmonic convolution. For the lower bound, clean alternatives with small prediction risk are transformed into legally edited null words whose restricted observations match those alternatives. The clean guarantee is used at the same horizon and supplied certificate as the attacked guarantee. Complete arguments appear in the supplement; the main text isolates the experiment and the information comparisons.

\section{Experiment and benchmark}\label{sec:experiment}
All logarithms are natural. Let $M\ge64$ be an integer and
\[
 X=(X_0,\ldots,X_{M+1}),\qquad Y_i=(X_i,X_{i+1}),\quad 0\le i\le M.
\]
Thus there are $M+2$ raw symbols, $M$ departures, and the terminal pair $Y_M$. A homogeneous order-two transition table $\theta_{ab}(c)$ induces the four-state pair kernel
$P((a,b),(b,c))=\theta_{ab}(c)$. There is a qualifying stationary law $\Pi$ such that
\begin{equation}\label{eq:certificate}
 \Pi P=\Pi,\qquad P^2(u,v)\ge\beta\Pi(v)\quad\text{for all }u,v\in\{00,01,10,11\}.
\end{equation}
The condition includes rows outside stationary support. Throughout the main results the supplied certificate is \emph{exactly} $\beta=1/M$; it need not be the source's largest possible coefficient. Neither $\theta$ nor $\Pi$ is revealed. Stationary initialization means $Y_0\sim\Pi$; arbitrary initialization means any hidden law $\mu$, including laws supported outside the stationary support.

\paragraph{Strict edits and actual-terminal loss.}
For $0\le k\le M$, a legal randomized kernel $K_{M,k}$ maps full binary words to full binary words. On every input $x$, every output $z$ in its support satisfies
\begin{equation}\label{eq:attack}
 z_M=x_M,\qquad z_{M+1}=x_{M+1},\qquad
 \sum_{i=0}^{M-1}\ind\{z_i\ne x_i\}\le k.
\end{equation}
The kernel is one fixed map on words, not supplied the true source table, stationary or initial law, or predictor's random seed. It may inspect the whole input and use public parameters and $k$; lower-bound kernels may hardcode a predeclared auxiliary family. No causal or computational restriction is imposed. Its randomness is independent of the predictor's. This formulation specifies the common kernels used in the proofs; source independence alone is not asserted to distinguish all standard contamination models.

For $Z\sim K_{M,k}(X,\cdot)$ and an estimator $A$ based on observation $\cT(Z)$, define
\begin{equation}\label{eq:risk}
 R(A;\theta,\mu,K_{M,k})
 =\E\,\TV\bigl(A(\cT(Z)),\theta_{Y_M}\bigr).
\end{equation}
The expectation includes the path and both independent randomizations. The target is the source's next-symbol law at the \emph{actual protected pair}, not an independent query. For binary distributions, TV is absolute error in either coordinate. Risk is unconditional.

\paragraph{Observed information.}
For $a,b,c\in\{0,1\}$ let
\[
 H_{abc}(z)=\sum_{i=0}^{M-1}\ind\{(z_i,z_{i+1},z_{i+2})=(a,b,c)\}.
\]
For prescribed $2\le r\le M+2$, the observation is
\begin{equation}\label{eq:observation}
 \cT_r(z)=\bigl((H_{abc}(z))_{a,b,c},\ z_0,\ldots,z_{r-1},\
 z_{M+2-r},\ldots,z_{M+1}\bigr).
\end{equation}
The windows may overlap. Counts and raw symbols are exact; neither edit locations nor the clean interior are supplied. In row $u=(a,b)$ use the designated successor $d(u)=1-a$, and define
\[
 J'_u=H_{ab0}(Z)+H_{ab1}(Z),\qquad C'_u=H_{ab,d(u)}(Z),\qquad
 q_u=\theta_{ab}(d(u)).
\]
Primes indicate corrupted counts. Unprimed counts are clean analysis variables. At the selected row $u=Y_M$, subscripts are omitted. A ratio $C/J$ is evaluated only when $J>0$.

\paragraph{One family for all budgets.}
Use the explicitly normalized rate profile
\begin{equation}\label{eq:benchmark}
 g_M(k)=\max\{M^{-1/2},k/M\},\qquad 0\le k\le M.
\end{equation}
For a prescribed sequence $r_\bullet=(r_M)_{M\ge64}$, let $C_{\rm stat}(r_\bullet)$ be the infimum of constants $C\ge0$ for which one estimator family $(A_M)$, based on $\cT_{r_M}$, satisfies
\begin{equation}\label{eq:factor}
 R(A_M;\theta,\Pi,K_{M,k})\le Cg_M(k)
\end{equation}
for every $M\ge64$, every $0\le k\le M$, every qualifying source and every legal kernel. Define $C_{\rm all}$ by additionally quantifying over every $\mu$. The family may depend on the prescribed sequence, $M$, $r_M$, and supplied $1/M$, but not on $k$ or hidden source quantities. The infimum need not be attained; an empty admissible set has value $+\infty$.

\paragraph{Finite-horizon competitive factors.}
Separately from the sequence factors $C$, fix $(M,r)$ and a single
$k$-blind map $A$ based on $\cT_r$. Define
\begin{equation}\label{eq:finite-factor}
 \Gamma_{\rm stat}(A;M,r)
 =\sup_{\substack{0\le k\le M,\ \theta,\Pi,K_{M,k}\\
                  \text{qualifying stationary experiments}}}
 \frac{R(A;\theta,\Pi,K_{M,k})}{g_M(k)},\qquad
 \Gamma_{\rm stat}(M,r)=\inf_A\Gamma_{\rm stat}(A;M,r).
\end{equation}
Define $\Gamma_{\rm all}(A;M,r)$ and $\Gamma_{\rm all}(M,r)$ by also
allowing every initial law $\mu$. The map may depend on $M,r$ and the
supplied $1/M$, but the very same map serves all $k$ at that horizon.
These are finite-horizon infima, not the sequence quantities $C$ above;
no infimum is assumed attained. Every map has factor at most $\sqrt M$
because TV is at most one.

The comparison is to a genuine same-information known-budget scale, rather than to an arbitrary denominator. Write $R^{\cT}_{\rm stat}(M,k)$ and $R^{\cT}_{\rm all}(M,k)$ for the respective minimax risks when $k$ is also supplied.
\begin{proposition}[Known-budget scale]\label{prop:oracle}
For the full corrupted history, for $\cT_r$, and for counts plus the protected pair,
\begin{equation}\label{eq:oracle}
 \frac{g_M(k)}{96}\le R^{\cT}_{\rm stat}(M,k)
 \le R^{\cT}_{\rm all}(M,k)\le191488g_M(k).
\end{equation}
\end{proposition}
The upper is count-only and allows arbitrary starts; the lower holds even for full histories. Thus a factor $C$ relative to $g$ gives a factor at most $96C$ relative to the corresponding minimax oracle. Conversely, a factor $D$ relative to that oracle gives a factor at most $191488D$ relative to $g$. Proofs are in Lemma~\ref{lem:oracle-floor} and Proposition~\ref{prop:count-oracle}.

For comparison with a more general certificate-dependent profile, define $F_2(M,0,\beta)=0$ and, for $k\ge1$,
\[
 F_2(M,k,\beta)=\frac{k}{M}\left[1+\log\left(1+
 \frac1{k\max\{\sqrt\beta,1/M\}}\right)\right],\qquad
 H_2=\min\{1,\max\{M^{-1/2},F_2\}\}.
\]
At the certificate used here, $\log(1+x)\le x$ gives $g_M(k)\le H_2(M,k,1/M)\le2g_M(k)$. Hence the two normalizations change optimal factors by at most two; all constants below are stated directly for $g$.

\section{Finite-horizon characterization and its consequences}\label{sec:results}
\begin{theorem}[Simultaneous finite-horizon window price]\label{thm:finite-window}
For every integer $M\ge64$ and $2\le r\le M+2$, put
$\ell=1+\log((M+2)/r)$. Then
\begin{equation}\label{eq:finite-window}
 2^{-30}\ell\le\Gamma_{\rm stat}(M,r)
 \le\Gamma_{\rm all}(M,r)\le2^{20}\ell.
\end{equation}
The upper uses only exact counts and the suffix. The lower grants both
boundary windows. All experiments use the supplied certificate $1/M$.
\end{theorem}
The proof in Section~\ref{sec:finite-window-proof} checks the complete finite
domain of the matched-observation bound. It tests the clean risk of the same
map before selecting an attacked budget; it is not a logarithmic risk lower
at each budget. Linear widths have constant factor; for fixed $0<\gamma<1$,
$r=\max\{2,\lceil M^\gamma\rceil\}$ has order $\log M$, as does $r=2$.
The constants are conservative, not a claim of sharp practical thresholds.

The next sequence statements are consequences at the order level. Their
original separate proofs are retained in
Proposition~\ref{prop:lower-consequences}, including the stronger
fixed-fraction lower constant.

\begin{theorem}[Exact raw-window criterion]\label{thm:width}
For every prescribed sequence of integers $2\le r_M\le M+2$,
\[
 C_{\rm stat}(r_\bullet)<\infty
 \quad\Longleftrightarrow\quad C_{\rm all}(r_\bullet)<\infty
 \quad\Longleftrightarrow\quad
 \liminf_{M\to\infty}\frac{r_M}{M}>0.
\]
No monotonicity of $r_M$ is required. The factor may depend on the prescribed sequence.
\end{theorem}
Necessity already uses stationary witnesses. Sufficiency includes unsupported starts. A finite initial set of horizons is harmless because $g_M(k)\ge M^{-1/2}$. The conclusion is about one budget-blind family, not a separate minimax estimator chosen for each $k$.

\begin{theorem}[Price of a fixed fraction]\label{thm:fraction}
Fix $0<\alpha\le1/4096$ and $r_M=\max\{2,\lceil\alpha M\rceil\}$. Abbreviating the corresponding factors by $C_{\rm stat}(\alpha)$ and $C_{\rm all}(\alpha)$,
\begin{equation}\label{eq:fraction}
 2^{-26}\log\frac1{32\alpha}
 \le C_{\rm stat}(\alpha)\le C_{\rm all}(\alpha)
 \le2^{19}\left(1+\log\frac2\alpha\right).
\end{equation}
The upper requires only the histogram and suffix; the lower permits both boundary windows.
\end{theorem}
Thus the optimal uniform factor has order $\log(1/\alpha)$ as $\alpha\downarrow0$. The upper is uniform in every $M,k$ for each fixed $\alpha$. The lower first sends $M\to\infty$ with $\alpha$ fixed and with a proposed finite factor fixed. This fixed-fraction proof takes limits in the stated order. The finite characterization above requires no limit in either parameter; neither result is a logarithmic risk lower bound at every budget. The constants are explicit but not optimized. Proposition~\ref{prop:lower-consequences} proves Theorems~\ref{thm:width} and~\ref{thm:fraction} from a finite upper bound and a finite matched-observation lower bound.

\subsection{An explicit estimator}
Set
\begin{equation}\label{eq:suffix-length}
 m=\min\{r-2,\lfloor M/2\rfloor\},\qquad N=M-m,\qquad
 L=1+\log\frac{M}{m+1}.
\end{equation}
The last $m+2$ symbols reveal exactly the $m$ departures indexed $M-m,\ldots,M-1$ and the terminal pair. Let $T'_u$ count visits to $u$ at those departures and put $O'_u=J'_u-T'_u$. Then $O'_u$ is the count on indices $0,\ldots,N-1$; it is a direct old-prefix count even though it is computed by subtraction. At the selected row, predict designated probability
\begin{equation}\label{eq:gate}
 \widehat q=
 \begin{cases}
 C'/J',& J'>0\ \text{and}\ 512O'\ge J',\\
 0,&\text{otherwise}.
 \end{cases}
\end{equation}
The complementary symbol receives probability $1-\widehat q$. This map does not use $k$ or $\beta$. When $r=2$, $m=0$ and $O'=J'$; the terminal pair is not a departure. Using $r-1$ in place of $r-2$ would overcount the available information.

The finite upper bound at the main certificate, proved in Proposition~\ref{prop:finite-upper}, is
\begin{equation}\label{eq:finite-upper}
 R\le\min\left\{1,\frac{131072}{\sqrt M}
 +(8448L+147744)\frac{k}{M}\right\}
 \le287264L g_M(k).
\end{equation}
If $r_M/M$ is eventually bounded below, then $L$ is eventually bounded. For fixed $\alpha$ as in Theorem~\ref{thm:fraction}, $m+1\ge\alpha M/2$, including the horizons at which $r_M=2$. This proves the two stated upper conclusions without a normalization loss.

\section{Why positions enter the proof}\label{sec:mechanisms}
\paragraph{Older visits and a truncated harmonic sum.}
The gate is not an estimate of the budget. It requires a positive fraction of the selected row's visits to precede the observed suffix, not merely one old occurrence. This proportional balance makes gate-switch errors chargeable to the edit budget; one forged old occurrence alone must not unlock a large corrupted suffix count. For a clean $d$-state homogeneous chain and any start, the relevant occupation inequality is
\begin{equation}\label{eq:reciprocal-main}
 \E\frac{\ind\{O\ge cJ\}}{J+1}
 \le(3d^2+2d^3)\frac{1+\log(M/(m+1))}{M},\qquad c>0.
\end{equation}
Here $J$ is the number of departures from the actual terminal state and $O$ counts those before time $M-m$. Lemma~\ref{lem:old-inverse} proves this by an exact first-hit convolution and a reciprocal-return identity. In this analytic bound the event can be enlarged to the existence of an old visit. That enlargement removes short backward-time terms; the surviving harmonic range produces the logarithm. The implemented gate still requires proportional balance. Generic finite-state first-hit pointmass estimates have precedents in \citet{NorrisPeresZhai2014}; the required gated convolution is proved explicitly.

A strict $k$-edit outcome changes each relevant count by at most $6k$. A deterministic true-gate/false-gate split therefore reduces the error to a clean empirical comparator, the reciprocal count in \eqref{eq:reciprocal-main}, a weighted imbalance term, and two small-occupation events. The all-state certificate controls the imbalance even from unsupported starts. No conditional independence of departures given the random terminal row or its occupation is assumed. The supplement's row-stack argument instead starts from independent row-wise randomness and then bounds the selected empirical deviation.

\paragraph{Mixture domination and legal observation matching.}
The lower bound uses stationary alternatives whose designated exit probabilities are
\[
 (q_{00},q_{01},q_{10},q_{11})=(a,1/2,1/2,t),\qquad
 \Pi_a=\frac{(t/a,2t,2t,1)}{1+t/a+4t}.
\]
An integral over $a$ creates the logarithmic term. Conditioning on a fixed departure count, one counts marked excursions and holding slots; the resulting occupation-mixture bound controls each proposed output mass. An interior loop relocation preserves the alternative's triple histogram and both of its observed windows, while making it close to a simple two-run null word. These are the alternative's windows, not windows of the unedited null input: changing an observed input-window symbol is legal but must be charged to the budget.

The attack normalizes only nonzero proposal masses, samples once, performs the relocation, and checks the actual Hamming cost and terminal pair. If any check fails, it returns its input unchanged. It is therefore a common source-independent kernel satisfying the budget on every outcome, not only on a high-probability event. In the risk comparison, nonnegative loss is discarded before proposal-mass domination; a signed triangle inequality is used only afterwards. Lemmas~\ref{lem:mixture} and~\ref{lem:relocation} and Proposition~\ref{prop:finite-lower} supply these steps.

For sublinear widths along any subsequence, a budget between the window and horizon scales makes the clean risk negligible relative to the alternative's exit scale while the mixture logarithm diverges. For fixed $\alpha$, a budget proportional to the window yields $\log(1/\alpha)$. The alternatives may depend on this tested budget, but their clean risk is bounded using the same estimator at the same $M$ and supplied $1/M$. This constrained-risk logic is related to general risk-transfer arguments such as \citet{DuchiRuan2021}; the dependent-word construction and legal observation matching are additional obligations.

\section{One computed positional bit}\label{sec:bit}
Now change the access model. The decoder receives exact triple counts and the protected pair, optionally also the initial pair. An encoder may scan the entire already corrupted word and send one symbol in a fixed finite alphabet $\mathcal A$. It knows $M$ but not $k$, $\beta$, or the source. The message is exact and fault-free after construction. There are no timing, framing, variable-length, or other side channels, and no shared source-dependent information.

\begin{corollary}[Minimum supplemental alphabet]\label{cor:onebit}
In this scan-and-message model, the minimum $|\mathcal A|$ for a finite uniform factor relative to $g_M(k)$ at supplied $1/M$ is two. This holds for stationary or arbitrary initialization, with or without the initial pair supplied.
\end{corollary}
For the upper, put $n=\lfloor M/2\rfloor$, select $u=Y_M$, and send
\begin{equation}\label{eq:bit}
 B=\ind\{512J'_1\ge J'\},\qquad
 J'_1=\sum_{i=0}^{n-1}\ind\{(Z_i,Z_{i+1})=u\}.
\end{equation}
The decoder outputs $C'/J'$ if $J'>0$ and $B=1$, and zero otherwise. Equality is accepted. If $J'=0$, the bit equals one, but the explicit positive-denominator test enforces the fallback. Proposition~\ref{prop:finite-upper} gives all-start risk at most
$\min\{1,2^{17}M^{-1/2}+2^{23}k/M\}$.
A one-symbol alphabet supplies no information. Theorem~\ref{thm:width} with $r_M=2$ rules it out even when both endpoint pairs are present. This proves minimality without another lower-bound construction.

\paragraph{A finite word illustration, not a risk experiment.}
For $M=65$, the two words in Table~\ref{tab:words} have the same initial pair $01$, terminal pair $10$, and triple counts, in lexicographic order,
\[
 (H_{000},H_{001},H_{010},H_{011},H_{100},H_{101},H_{110},H_{111})
 =(0,1,1,1,1,0,1,60).
\]
At row $10$, $J'=C'=1$. The single departure from that row lies in different halves, so the computed gate differs although $\cT_2$ is identical. The example verifies a combinatorial collision only: it gives no source probabilities, statistical risk comparison, or lower-bound proof.
\begin{table}[ht]
\centering
\caption{An exact $67$-symbol collision. The notation $1^{62}$ means $62$ consecutive ones; the first half contains departures $0$ through $31$.}\label{tab:words}
\begin{tabular}{lccc}
\toprule
Word & $J'_1$ & $B$ & Designated prediction\\
\midrule
$0100\,1^{62}\,0$ & $1$ & $1$ & $1$\\
$0\,1^{62}\,0010$ & $0$ & $0$ & $0$\\
\bottomrule
\end{tabular}
\end{table}

This is a worst-case supplemental alphabet result, not an entropy or mutual-information statement. The baseline exact counts are uncharged inputs. One computed bit is not one observed raw coordinate and does not remove the whole-word scan. A counter-based implementation is a natural baseline: maintain the eight triple counts and four first-half row counts, retain the endpoint, and select its row after the scan. This uses $O(\log M)$ counter workspace, $O(M)$ counter updates, and at most $O(M\log M)$ naive bit operations. No neural mechanism is needed for this statistical rate guarantee.

\section{Relation to prior work and scope}\label{sec:related}
\paragraph{Clean positional prediction.}
\citet[model (1), Theorem 24, and Section VI.A]{HanJanaWu2023} study stationary prediction at the actual terminal state, including suffix-based estimators and higher-order KL analysis. Theorem 24 gives the binary order-two KL result; the squared-TV result in Section VI.A concerns first-order chains with fixed state count at least three. These are direct antecedents, not independent-query results. Their clean full-trajectory experiment differs from the strict-edit and restricted-observation comparison here. Losses must also be distinguished: a squared-TV upper bound implies an unsquared-TV upper bound by Jensen, but a squared-TV lower bound does not yield its square root as an unsquared-TV lower bound. The spectral-gap-restricted results in that work use an assumption different from the supplied all-state minorization \eqref{eq:certificate}.

\paragraph{Unknown contamination.}
Budget adaptation is not new in itself. \citet{JainOrlitskyRavindrakumar2022} select among candidate estimates with monotone radii under a simultaneous-validity premise: every valid-index ball must contain the target. Separate expected-risk bounds do not establish that premise. Their selection lemma is pointwise, however, so the random terminal target is not by itself an obstruction: it could be applied realization by realization if all candidate predictions share the protected query and simultaneous accuracy is proved. \citet{BhattEtAl2022} explicitly construct confidence-based adaptation for a mean under strict retrospective replacements of IID observations. \citet{ChenGaoRen2016} give robust testing and estimation under IID Huber mixtures, while \citet{LiuGao2019} exhibit an adaptation obstruction involving both smoothness and contamination. None of these comparisons licenses a general claim that unknown contamination alone prevents adaptation.

\paragraph{Dependence and count-based inference.}
\citet{Lecestre2023} estimate marginal or invariant distributions under dependence; their contamination corollary imposes independence conditions on the replacement channel. \citet{LeskelaDreveton2026} estimate transition matrices from independent parallel trajectories with arbitrarily corrupted whole trajectories. These are not the one-word terminal-target experiment. They also show why the present result should not be summarized as a universal failure of counts under corruption.

Clean transition-count factorization and Markov exchangeability remain valid. The report of \citet{DiaconisFreedman1978} and the finite-stack discussion of \citet{HalberstamHutchcroft2024} provide relevant background on count-compatible paths. For a clean order-two word, likelihood is its initial-pair mass times transition probabilities raised to triple counts. After a position-sensitive kernel $K$, an observed word's probability instead sums these clean probabilities weighted by $K(z\mid x)$. Clean factorization alone does not establish sufficiency in this altered experiment. The specific observation map, not an unrestricted claim about all compact summaries, is essential.

\paragraph{Limits of the conclusions.}
The lower bounds concern uniform unknown-budget comparison at the critical supplied certificate, with exact counts and fixed boundary windows. They are not full-history impossibility results, every-budget logarithmic minimax penalties, claims about arbitrary interior statistics, or noisy-bit guarantees. The positive results include arbitrary starts; the lower witnesses are stationary. The constants are deliberately conservative. The word example is not empirical validation, and no learned-predictor experiment is asserted.

Computational realizations may be studied separately, but resource distinctions do not change these statistical statements. A whole-word counter baseline must be compared on the same input and arithmetic model. A finite-precision network realization would additionally need to account for width, internal precision, scanning, storage, and decoder cost. Neither existence of such a realization nor a depth bound would establish learnability, optimization success, Transformer superiority, or a general architecture lower bound. The present manuscript makes none of those claims.

\clearpage
\appendix
% MANUSCRIPT DRAFT: pending complete independent consolidated review.
% This file contains no external proof dependencies on project-private units.
\section{Finite-state estimates and upper bounds}
\label{sec:positive-proofs}

Throughout, a departure count excludes the terminal state.  For a chain
$(Y_i)_{i=0}^M$, write
\[
 J_u=\sum_{i=0}^{M-1}\mathbf1\{Y_i=u\},\qquad J=J_{Y_M}.
\]
For the binary order-two experiment, $Y_i=(X_i,X_{i+1})$,
$X_0,\ldots,X_{M+1}$ contains $M$ departures, and
$q_u=\theta_u(d(u))$ with $d(a,b)=1-a$.  Let $C_u$ count the
corresponding designated successors and put $C=C_{Y_M}$ and $q=q_{Y_M}$.
Primed counts are computed on the corrupted word.  Every bound below permits
an arbitrary initial law.  The all-state certificate, when invoked, is
$P^2(u,v)\ge\beta\Pi(v)$ for every pair of states, with $\beta>0$.
The final pair is protected, so clean and corrupted counts always select the
same actual query row.

\subsection{Universal occupation and clean prediction}

\begin{lemma}[Actual-terminal occupation]\label{lem:occupation}
For any homogeneous chain on $d$ states, any initial law, any $M\ge0$, and
any integer $w\ge1$,
\[
 \Pr(J_{Y_M}<w)\le\min\left\{1,\frac{d^2w}{M+1}\right\}.
\]
No irreducibility or stationary-start assumption is needed.
\end{lemma}
\begin{proof}
For each $t$ set $J(t)=\sum_{i<t}\mathbf1\{Y_i=Y_t\}$,
$b_t(v)=\Pr_v(J(t)<w)$, and $B_t=\max_v b_t(v)$.  If $t\ge\ell$,
the full count dominates the count in the final $\ell$-step suffix.
Conditioning on the suffix's initial state and using homogeneity gives
$B_t\le B_\ell$.  On every path, each state can occur at most $w$ times
with fewer than $w$ earlier occurrences.  Hence
$\sum_{t=0}^M\mathbf1\{J(t)<w\}\le dw$.  Consequently
\[
 (M+1)B_M\le\sum_{t=0}^M B_t
 \le\sum_v\sum_{t=0}^M b_t(v)\le d^2w.
\]
Average over the initial law and use the probability cap.
\end{proof}

\begin{lemma}[Binary clean empirical comparison]\label{lem:empirical}
For every homogeneous binary order-two source, every initial law, and
$M\ge64$,
\[
 \E\left[\mathbf1\{J>0\}\left|\frac CJ-q\right|\right]
 \le\frac{512}{\sqrt{M+1}}.
\]
The expression inside the indicator is evaluated only when $J>0$.
\end{lemma}
\begin{proof}
For each of the four rows $u$, generate an independent stack of $2M$
Bernoulli$(q_u)$ variables $\xi_{u,j}$, independently of the initial pair.
At the $j$th departure from $u$, use $\xi_{u,j}$ to choose whether the
successor is $d(u)$.  For every fixed finite path, the used stack prefixes
have exactly the product of its transition probabilities; thus this realizes
the original law, including boundary probabilities.  Write
$S_{u,j}=\sum_{i=1}^j(\xi_{u,i}-q_u)$.  Pathwise,
$C-Jq=S_{Y_M,J}$ on $J>0$.  The stacks are not asserted independent after
conditioning on $J$, $Y_M$, or any selection event.

For a centered variable in an interval of length one, its log moment
generating function has second derivative equal to a tilted variance,
which is at most $1/4$.  Integrating twice gives
$\E e^{\lambda V}\le e^{\lambda^2/8}$.  Thus
$e^{\lambda S_{u,j}-j\lambda^2/8}$ is a nonnegative supermartingale.
Stopping at its first level crossing before the deterministic time $2v$
gives, for $x>0$,
\[
 \Pr\left(\max_{j\le2v}S_{u,j}\ge x\right)
 \le e^{-\lambda x+2v\lambda^2/8}
 \le e^{-x^2/v},\qquad \lambda=2x/v.
\]
The stopping inequality here is finite: the expectation of the stopped
process is at most its initial value by summing conditional expected
increments through time $2v$.  Apply the same argument to $-S$ and take
a union over rows.  For $Z_v=\max_u\max_{j\le2v}|S_{u,j}|$,
\begin{equation}\label{eq:p-stack-tail}
 \Pr(Z_v\ge x)\le8e^{-x^2/v}.
\end{equation}
For dyadic integers $v\le M$, let $E_v=\{v\le J<2v\}$ and
$p_v=\min\{1,32v/(M+1)\}$.  Lemma~\ref{lem:occupation} gives
$\Pr(E_v)\le p_v$.  Without any independence assumption, tail integration
and $a=\sqrt{v\log(8/p_v)}$ yield
\[
 \E[Z_v\mathbf1_{E_v}]
 \le\int_0^\infty\min\{p_v,8e^{-x^2/v}\}\,dx
 \le p_va+\frac{p_vv}{2a}
 \le2p_v\sqrt{v\log(8/p_v)}.
\]
For the middle inequality, bound the Gaussian tail above $a$ by inserting
$x/a\ge1$ and integrating.  This layer contributes at most
$2p_v\sqrt{\log(8/p_v)}/\sqrt v$ to the desired risk.
Put $N=M+1$.  Summing layers with $v\ge N/32$ costs at most
\[
 \frac{2\sqrt{32\log8}}{1-2^{-1/2}}N^{-1/2}<56N^{-1/2}.
\]
If $v_0$ is the largest dyadic integer below $N/32$, then
$N/64\le v_0<N/32$ and
$\log(N/(4v_02^{-j}))\le3(j+1)$.  The remaining layers cost at most
\[
 \frac{64\sqrt3}{\sqrt{32N}}
 \sum_{j\ge0}2^{-j/2}\sqrt{j+1}
 \le\frac{64\sqrt3}{\sqrt{32N}}(1-2^{-1/2})^{-2}
 <236N^{-1/2}.
\]
Only a numerical series was extended; all stack maxima use at most $2M$
variables.  The layers partition $1\le J\le M$, including the last layer
when $M$ is dyadic.  The safe constant $512$ follows.
\end{proof}

\subsection{Return cycles and binary hitting scales}

\begin{lemma}[Finite-state return identity]\label{lem:kac}
In a finite irreducible chain with stationary law $p$, the first positive
return $T$ to $u$ has mean $\E_uT=1/p_u$.  If
$H_u=\max_v\E_vD_u$, where $D_u$ is the ordinary hitting time, zero at
$u$, then $1/p_u\le1+H_u$.
\end{lemma}
\begin{proof}
Irreducibility and finiteness give a fixed block length and a positive
minimum probability of hitting $u$ within a block, uniformly over starting
states.  Repeated conditioning gives geometric tails for hitting and return
times, so the following expectations are finite.  Let
$\nu(v)=\E_u\sum_{j=0}^{T-1}\mathbf1\{Y_j=v\}$.  First-step conditioning
inside this sum gives
\[
 (\nu P)(v)=\E_u\sum_{j=1}^{T}\mathbf1\{Y_j=v\}=\nu(v),
\]
since $Y_0=Y_T=u$.  A finite irreducible chain has a unique stationary
probability: for completeness, if $\eta$ is another one, the ratios
$\eta(v)/p_v$ are harmonic for the reversed transition matrix; a maximal
ratio propagates along every positive reversed edge and hence is constant.
Thus $\nu=(\E_uT)p$.  Exactly one visit to $u$ occurs before $T$, so
$\nu(u)=1$, proving the identity.  A positive return is one step followed
by an ordinary hit, proving $\E_uT\le1+H_u$.
\end{proof}

\begin{lemma}[Support and weighted hitting scale]\label{lem:hitting-scale}
A binary order-two source with a positive all-state certificate has either
full stationary support or a singleton absorbing diagonal support.  In the
full-support case, with $q=(a,b,e,\tau)$ and
$\overline H_u=\max\{H_u,1/\Pi(u)\}$,
\begin{equation}\label{eq:p-weighted-hitting}
 \sum_u\Pi(u)q_u\overline H_u\le\frac{40}{\sqrt\beta}.
\end{equation}
\end{lemma}
\begin{proof}
Stationarity of consecutive pair marginals gives
$\Pi=(y,x,x,w)$, and stationary flow into each diagonal gives
\begin{equation}\label{eq:p-flow}
 ya=xe,\qquad w\tau=xb.
\end{equation}
If $x>0$, the certificate entries
$P^2(10,01)=ea\ge\beta x$ and
$P^2(01,10)=b\tau\ge\beta x$ force $a,b,e,\tau>0$.
The flow identities then imply $y,w>0$.  If $x=0$, the stationary law
is on diagonals; each supported diagonal is absorbing by stationarity.
Two absorbing supported diagonals contradict the positive certificate.
This proves the classification.  Summing a certificate row gives
$\beta\le1$.  Full support and the certificate imply irreducibility.

Assume full support and put $\gamma=(a,1,1,\tau)\ge q$.
The two certificate entries just used, together with
\eqref{eq:p-flow}, also give
\begin{equation}\label{eq:p-square-entry}
 e^2\ge\beta y,\qquad b^2\ge\beta w.
\end{equation}
Here, for example, $ea\ge\beta x=\beta ya/e$ and $a,e>0$.
For target $00$, write $h_{01},h_{10},h_{11}$ for the mean hit times.
The first-step equations are
\[
 h_{01}=1+(1-b)h_{10}+bh_{11},\quad
 h_{10}=1+(1-e)h_{01},\quad h_{11}=\tau^{-1}+h_{10}.
\]
They give
\[
 h_{01}=\frac{b+2\tau}{e\tau},\quad
 h_{10}=h_{01}-\frac b\tau-1,\quad
 h_{11}=h_{01}+\frac{1-b}{\tau}-1,
\]
and therefore $H_{00}\le b/(e\tau)+2/e+1/\tau+1$.
For target $01$, the three hit times from $00,10,11$ are
$1/a$, $1+e/a$, and $1/\tau+1+e/a$, so
$H_{01}\le2/a+1/\tau+2$.
Exchanging symbols gives the other two bounds:
\[
 H_{10}\le1/a+2/\tau+2,\qquad
 H_{11}\le e/(ba)+2/b+1/a+1.
\]
The substitution exchanges $00$ with $11$, $01$ with $10$, and
$(a,b,e,\tau)$ with $(\tau,e,b,a)$; it does not assume $b=e$.
Using \eqref{eq:p-flow} and \eqref{eq:p-square-entry},
\begin{align*}
 \sqrt\beta\,yaH_{00}
 &\le\sqrt\beta\,(w+2x+ew/b+xe)
 \le\sqrt\beta\,(w+3x)+\sqrt w\le5,\\
 \sqrt\beta\,xH_{01}
 &\le2\sqrt\beta\,y/e+\sqrt\beta\,w/b+2\sqrt\beta\,x
 \le2\sqrt y+\sqrt w+2x\le5.
\end{align*}
The symbol-exchanged inequalities give the same bound for $10,11$.
Lemma~\ref{lem:kac} implies $\overline H_u\le H_u+1$; hence
$\Pi(u)\gamma_u\overline H_u\le6/\sqrt\beta\le10/\sqrt\beta$.
Sum over four states and use $q\le\gamma$.
\end{proof}

\subsection{First hits, endpoint probabilities, and joint tails}

\begin{lemma}[Target endpoint and occupation bounds]\label{lem:endpoint-tail}
For any $d$-state chain, possibly reducible, and $v\ne u$,
\begin{equation}\label{eq:p-first-hit}
 \Pr_v(D_u=t)\le\frac d{t+1},\qquad t\ge1.
\end{equation}
For an irreducible chain, fix a target with stationary mass $p>0$,
maximum ordinary hitting mean $H$, and $h=\max\{H,1/p\}$.
For every initial state $i$ and every integer $\ell\ge1$, with
$J_u(\ell)=\sum_{j=0}^{\ell-1}\mathbf1\{Y_j=u\}$,
\begin{align}
 \Pr_i(Y_\ell=u)&\le4dp(1+h/\ell),\label{eq:p-endpoint}\\
 \Pr_i(J_u(\ell)<\ell p/16)&\le2e^{-\ell/(8h)},\label{eq:p-lower-tail}\\
 \Pr_i(Y_\ell=u,J_u(\ell)>4\ell p+1)
 &\le16dp(1+h/\ell)e^{-\ell/(2h)}.\label{eq:p-joint-tail}
\end{align}
\end{lemma}
\begin{proof}
\emph{Reciprocal return coefficients and first hits.}
For any finite stochastic matrix $P$, let $T_j$ be the $j$th positive
return to $u$, allowing defective returns, and let
$F(z)=\E_u[z^{T_1};T_1<\infty]$.  Strong Markov at returns gives
$\Pr_u(T_j=s)=[z^s]F(z)^j$.  All strong-Markov uses here can be obtained
by conditioning on each possible stopping time and state and multiplying
the subsequent transition probabilities; no conditional independent-sample
claim is involved.  Let $P_{-u}$ delete row and column $u$.  Writing $P$
in target/rest block form, direct summation over the excursion before the
first return gives
$F(z)=zP_{uu}+z^2P_{u,-u}(I-zP_{-u})^{-1}P_{-u,u}$.
The block determinant identity therefore yields
\[
 1-F(z)=\frac{\det(I-zP)}{\det(I-zP_{-u})}.
\]
These are formal series, so no convergence issue arises.  The identity
$-\log\det(I-zA)=\sum_{s\ge1}\operatorname{tr}(A^s)z^s/s$
follows by differentiating the determinant, expanding $(I-zA)^{-1}$,
and integrating from zero.  Consequently
\begin{equation}\label{eq:p-return-trace}
 \sum_{j\ge1}\frac{\Pr_u(T_j=s)}j
 =\frac{\operatorname{tr}(P^s)-\operatorname{tr}(P_{-u}^s)}s
 \le\frac ds,\qquad s\ge1.
\end{equation}
Its left side is nonnegative; the upper bound uses nonnegative entries
of $P_{-u}^s$ and diagonal entries of $P^s$ at most one.
Now change only row $u$ to send $u$ deterministically to $v\ne u$.
In this auxiliary stochastic chain, the first positive return has law
$1+D_u$ under the original start $v$.  Retain only the first-return
term of \eqref{eq:p-return-trace} to prove \eqref{eq:p-first-hit}.
The auxiliary chain need not be an order-two source.

\emph{A target-specific endpoint bound.}
Return now to the irreducible chain.  Let $g(v)=\E_vD_u$.
First-step equations and Lemma~\ref{lem:kac} give
$(Pg-g)(v)=-1+\mathbf1\{v=u\}/p$.  Telescoping expectations shows
\begin{equation}\label{eq:p-expected-count}
 \E_vJ_u(s)=p\{s+\E_vg(Y_s)-g(v)\}\le p(s+H).
\end{equation}
A loop-rotation identity is
\[
 sP^s(u,u)=\sum_v\E_v[J_u(s)\mathbf1\{Y_s=v\}]
 \le dp(s+H).
\]
Indeed, for each $j<s$, the corresponding sum is
$\sum_vP^j(v,u)P^{s-j}(u,v)=P^s(u,u)$.
For $v\ne u$, split the first-hit convolution at $t=\lfloor s/2\rfloor$.
For early hits the return factor is at most $dp(1+2H/s)$.
For late hits \eqref{eq:p-first-hit} is at most $2d/s$, and
\eqref{eq:p-expected-count} from $u$ bounds the sum of return factors
by $p(\lceil s/2\rceil+H)$.  The resulting bound is at most
$dp(2+2/s+4H/s)\le4dp(1+H/s)$.  It also holds for $v=u$;
the early sum may be empty at $s=1$.  This proves the slightly stronger
version of \eqref{eq:p-endpoint} with $H$ instead of $h$.

\emph{Return moments and positive transforms.}
The family consisting of ordinary hitting times from nontarget states and
the positive return from $u$ is closed under conditioning on survival and
the current state.  Its means are at most $h$.  For each member $V$ and
integer $r\ge1$,
\begin{equation}\label{eq:p-return-moment}
 \E V^r\le r!h^{r-1}\E V.
\end{equation}
To prove this by induction, use
$V^r\le r\sum_{j=0}^{V-1}(V-j)^{r-1}$, condition each summand on
survival through time $j$, and apply the induction hypothesis and the
uniform mean bound $h$ to the residual law.  The sum of the survival
probabilities is $\E V$.
With $\lambda=1/(4h)$, expansion in nonnegative powers gives
\begin{equation}\label{eq:p-return-transforms}
 \E_i e^{\lambda D_u}\le\frac43,\qquad
 \E_u e^{\lambda T_1}\le e^{2\lambda/p},\qquad
 f:=\E_u e^{-\lambda T_1}\le e^{-3\lambda/(4p)}.
\end{equation}
For the last inequality use $e^{-x}\le1-x+x^2/2$ and
$\E T_1^2\le2h/p$; the first two follow from the geometric-series
bound on the moments.  If $K=\lceil\ell p/16\rceil$, fewer than $K$
visits in $0,\ldots,\ell-1$ means
$D_u+T_1+\cdots+T_{K-1}\ge\ell$.
These are fresh independent return cycles after the first hit, before any
horizon restriction.  Exponential Markov and \eqref{eq:p-return-transforms}
bound this probability by
$(4/3)e^{-\lambda\ell+2\lambda(K-1)/p}
\le(4/3)e^{-7\lambda\ell/8}\le2e^{-\ell/(8h)}$.
This includes $K=1$ and proves \eqref{eq:p-lower-tail}.

\emph{A joint tail by an exact change of path measure.}
Let $h_i=\E_i e^{-\lambda D_u}$, with $h_u=1$.
The inequality $e^{-x}\ge1-x$ gives $h_i\ge3/4$ and $f\ge3/4$.
Define
\[
 Q_{ij}=e^{-\lambda}P_{ij}h_j/h_i\quad(i\ne u),
 \qquad Q_{uj}=e^{-\lambda}P_{uj}h_j/f.
\]
The hitting and return equations show that every row sums to one.
This auxiliary matrix preserves directed support, so is irreducible.
Its ordinary hitting law is the original law tilted by
$e^{-\lambda D_u}/h_i$, whence $H_Q\le2H$.
Its positive return law is tilted by $e^{-\lambda T_1}/f$, and
\[
 \frac1{2p}\le\E_QT_1\le\frac2p.
\]
For the lower bound, use
$\E[T_1e^{-\lambda T_1}]\ge1/p-\lambda(2h/p)=1/(2p)$
and $f\le1$; the upper bound uses $f\ge3/4$.
Kac for $Q$ gives $p_Q\le2p$ and
$\max(H_Q,1/p_Q)\le2h$.  The endpoint bound already proved yields
$Q^\ell(i,u)\le16dp(1+h/\ell)$.
For every length-$\ell$ path from $i$ to $u$, telescoping the factors gives
\[
 Q(\text{path})=P(\text{path})e^{-\lambda\ell}
                  f^{-J_u(\ell)}/h_i.
\]
The count is exactly the number of target departures.  Thus on
$J_u(\ell)>4\ell p+1$, summing the reverse likelihood relation gives
\[
 \Pr_i(Y_\ell=u,J_u(\ell)>4\ell p+1)
 \le h_ie^{\lambda\ell}f^{4\ell p}Q^\ell(i,u)
 \le16dp(1+h/\ell)e^{-\ell/(2h)}.
\]
This proves the joint bound with its stationary-mass factor; no product
of an unconditional upper tail and an endpoint probability was taken.
\end{proof}

\subsection{Weighted imbalance, including unsupported starts}

\begin{lemma}[Half-window imbalance]\label{lem:imbalance}
For $M\ge64$, put $n=\lfloor M/2\rfloor$ and
$J_{1,u}=\sum_{i=0}^{n-1}\mathbf1\{Y_i=u\}$.
Every binary source with a positive all-state certificate, under every
initial law, satisfies
\[
 \E\left[q_{Y_M}\mathbf1\{J>0,J_{1,Y_M}<J/256\}\right]
 \le\frac{65536}{M\sqrt\beta}.
\]
\end{lemma}
\begin{proof}
We first prove a general irreducible-chain statement.  Let $m=M-n$,
$J_{2,u}=\sum_{i=n}^{M-1}\mathbf1\{Y_i=u\}$, and assign arbitrary
weights $0\le q_u\le1$.  Write
$W_u=p_uq_u\overline H_u$, where $p_u$ is stationary mass.
For every $0<c\le1/256$ we claim
\begin{equation}\label{eq:p-generic-imbalance}
 \E\left[q_{Y_M}\mathbf1\{J_{1,Y_M}<cJ\}\right]
 \le\frac{288d}{M}\sum_uW_u.
\end{equation}
The strict imbalance event is empty at $J=0$.
Fix $u$ with $np_u\ge1$.  If
$J_{1,u}\ge np_u/16$ and $J_{2,u}\le4mp_u+1$, then $m\le2n$ gives
$J_{1,u}/(J_{1,u}+J_{2,u})\ge1/145>1/256$.
Thus imbalance requires a first-window lower tail or a second-window
upper tail.  For the first event, condition on the whole history through
$Y_n$, use the uniform future endpoint bound
\eqref{eq:p-endpoint}, and then the first-window lower tail
\eqref{eq:p-lower-tail}.  For the second event condition only on $Y_n$
and use the joint bound \eqref{eq:p-joint-tail}.  Since $m\ge n$,
\[
 \Pr_i(Y_M=u,\ J_{1,u}<cJ_u)
 \le24dp_u(1+\overline H_u/n)e^{-n/(8\overline H_u)}.
\]
After multiplication by $q_u$, this is
$(24dW_u/n)(x+1)e^{-x/8}$, where $x=n/\overline H_u$.
Since the maximum of $(x+1)e^{-x/8}$ is
$8e^{-7/8}<4$, it is at most $96dW_u/n$.
If $np_u<1$, use only the endpoint bound at time $M$.
As $W_u\ge q_u$ and $p_uq_u<q_u/n$,
\[
 q_u\Pr_i(Y_M=u)\le4d(p_uq_u+W_u/M)\le8dW_u/n.
\]
Sum over $u$, average over initial states, and use $n\ge M/3$ to prove
\eqref{eq:p-generic-imbalance}.
For full-support binary sources, Lemma~\ref{lem:hitting-scale} now gives
$288\cdot4\cdot40/(M\sqrt\beta)=46080/(M\sqrt\beta)$.

It remains to treat singleton stationary support, without discarding
unsupported actual initial states.  Suppose $00$ is absorbing, so $a=0$,
and denote the other states by $A=01$, $B=10$, $C=11$.
An original path ending in this active set has never hit $00$.
Reflect the killing edge $B\to00$ to $B\to A$, producing the auxiliary
three-state chain
\[
 R(A,B)=1-b,\quad R(A,C)=b,\quad R(B,A)=1,
 \quad R(C,B)=\tau,\quad R(C,C)=1-\tau.
\]
Every original active path has probability at most its reflected-path
probability: the only changed factor is $1-e\le1$ on $B\to A$.
The endpoint and all active-state counts are unchanged.  Thus any
nonnegative functional of these paths is dominated without conditioning
on survival.
For $b,\tau>0$ this chain is irreducible, with
\begin{equation}\label{eq:p-reflected-weights}
 Z=2\tau+b,\quad p_A=p_B=\tau/Z,\quad p_C=b/Z,
 \qquad p_Ab\overline H_A\le1,\quad
 p_C\tau\overline H_C\le1.
\end{equation}
Indeed $H_A=1+1/\tau$ and $H_C=2/b$ by first-step equations;
$\overline H_A\le2+1/\tau$ and
$\overline H_C\le1+2/b$ give the last inequalities since $b,\tau\le1$.
Apply \eqref{eq:p-generic-imbalance} with weights $(b,0,\tau)$:
the $A,C$ contribution is at most $288\cdot3\cdot2/M=1728/M$.
If $b=0$ or $\tau=0$, approximate these two parameters by positive
ones in the reflected chain.  At fixed finite horizon, path probabilities
and the weights are continuous, and the bound is uniform, so the same
conclusion follows by taking a limit.  No minorization-preserving
approximation of the original source is being asserted.
The remaining target $B$ has weight $e$, and for an active start
\begin{equation}\label{eq:p-next-absorption}
 e\Pr(Y_M=B)=\Pr(D_{00}=M+1)\le\frac4{M+2},
\end{equation}
by \eqref{eq:p-first-hit}: the sole entry edge to $00$ is $B\to00$.
Drop its imbalance event.  An initial $00$ contributes zero; mixtures
follow by linearity.  The total is at most $1732/M$, below the stated
bound because $\beta\le1$.  Exchanging symbols handles absorbing $11$
with the same fixed designated-successor rule.
\end{proof}

\subsection{Reciprocal occupation with an old visit}

\begin{lemma}[Old-visit reciprocal count]\label{lem:old-inverse}
For any homogeneous $d$-state chain and initial law, let $M\ge2$,
$0\le m\le\lfloor M/2\rfloor$, and
\[
 O_u=\sum_{i=0}^{M-m-1}\mathbf1\{Y_i=u\},\qquad O=O_{Y_M}.
\]
For every $c>0$,
\begin{align*}
 \E\frac{\mathbf1\{O\ge cJ\}}{J+1}
 &\le\frac1M\left(3d^2+2d^3
     \sum_{s=m+1}^{\lfloor(M-1)/2\rfloor}\frac1s\right)\\
 &\le\frac{3d^2+2d^3}{M}
       \left(1+\log\frac M{m+1}\right),
\end{align*}
where an empty sum is zero.  In particular, the last numerator constant
is $176$ for the pair chain.  The statement includes $J=0$ and $m=0$.
\end{lemma}
\begin{proof}
Starting at $u$, the event
$\{Y_s=u,J_u(s)=j\}$ is exactly $\{T_j=s\}$ for $j\ge1$:
the initial visit is a departure, while the terminal visit is not.
Equation~\eqref{eq:p-return-trace}, which allowed defective returns, gives
\[
 K_u(s):=\E_u\frac{\mathbf1\{Y_s=u\}}{J_u(s)+1}
 =\sum_{j\ge1}\frac{\Pr_u(T_j=s)}{j+1}\le\frac ds.
\]
Fix an initial state $v\ne u$, and put
$a_v(t)=\Pr_v(D_u=t)$.  The event $Y_M=u,J_u=0$ means $D_u=M$,
so its contribution is at most $d/(M+1)\le d/M$ by
\eqref{eq:p-first-hit}; it is impossible if $v=u$.
On $J_u>0,O_u\ge cJ_u$, integrality forces $O_u\ge1$, hence the first
hit is at most $M-m-1$.  Discarding the rest of the gate and using
strong Markov gives the exact identity
\[
 \E_v\frac{\mathbf1\{Y_M=u,O_u\ge1\}}{J_u+1}
 =\sum_{t=1}^{M-m-1}a_v(t)K_u(M-t).
\]
There is no target departure before $t$, so the full count is exactly
the post-hit departure count, with no additional unit shift.
For $t\le M/2$, use $K_u(M-t)\le2d/M$ and
$\sum_ta_v(t)\le1$; the contribution is at most $2d/M$.
This also covers $v=u$, whose expression is just $K_u(M)$.
For $t>M/2$, use $a_v(t)\le2d/M$ and write $s=M-t$;
the contribution is at most
\[
 \frac{2d^2}{M}\sum_{s=m+1}^{\lfloor(M-1)/2\rfloor}\frac1s.
\]
Together with the zero-count term this is a per-target bound.
Average over the start and sum over $d$ terminal targets to obtain the
first inequality.  Finally,
$\sum_{s=a}^b1/s\le1/a+\log(b/a)\le1+\log(M/a)$ for $b\le M$,
with the empty case automatic, gives the stated logarithmic bound.
\end{proof}

\subsection{Protected edits and the two positional gates}

\begin{lemma}[Pathwise gate reduction]\label{lem:gate}
Let $D_u$ be a count of row $u$ over a fixed set of departure indices,
and at the protected row put $D=D_{Y_M}$.  On the corrupted word use
\[
 \widehat q=\begin{cases}
 C'/J',&J'>0\ \hbox{and }512D'\ge J',\\
 0,&\hbox{otherwise}.
 \end{cases}
\]
For every allowed corruption kernel and every initial law,
\begin{align}\label{eq:p-gate-reduction}
 R\le{}&\E\left[\mathbf1\{J>0\}|C/J-q|\right]
 +48k\E\frac{\mathbf1\{D\ge J/1024\}}{J+1}
 +\E[q\mathbf1\{J>0,D<J/256\}]\notag\\
 &+\Pr(J<6156k+1)+\Pr(J<3078k+1).
\end{align}
In particular this applies to a first-half count or to the old-prefix
count reconstructed by subtracting a suffix count from the total.
\end{lemma}
\begin{proof}
One edited symbol affects at most three departure triples; each changed
triple changes the histogram's $\ell^1$ norm by at most two.  Pair-count
windows affect at most two departures.  Taking a union over the edits,
including adjacent edits and indices straddling a split, gives the safe
bounds
\[
 |J'-J|,\ |C'-C|,\ |D'-D|\le\delta:=6k.
\]
When the old count is computed by subtraction, it is still exactly a
direct count on each word, so the same bound applies without adding
two separate perturbation bounds.  Endpoint protection keeps the selected
row fixed.
Let $G'=\{J'>0,512D'\ge J'\}$,
$H=\{D\ge J/1024\}$, and
$B=\{J>0,D<J/256\}$.  On $G'\cap H$ with $J>0$,
\[
 |C'/J'-q|\le|C/J-q|+\frac{8\delta}{J+1}.
\]
For $\delta=0$ the ratios coincide.  For $\delta>0$ and $J\ge2\delta$,
$J'\ge J/2$ and ratio subtraction gives a difference at most
$4\delta/J\le8\delta/(J+1)$.  For $J<2\delta$, use the unit bound on
the difference of two probabilities.
On $G'\cap H^c$,
$D+\delta\ge D'\ge(J-\delta)/512$ whereas $D<J/1024$,
so $J<1026\delta=6156k$.
On $(G')^c\cap B^c$ with $J>0$, if $J'=0$ then $J\le\delta$;
otherwise $D'<J'/512$ and $D\ge J/256$ give
$J/256-\delta<(J+\delta)/512$, hence $J<513\delta=3078k$.
If the gate rejects on $B$, its loss is exactly $q$.
The two displayed integer-threshold tails also cover $J=0$.
These cases give \eqref{eq:p-gate-reduction} pathwise and then in
expectation, with $8\delta=48k$.  At $k=0$ the false-status cases for
$J>0$ are impossible.  Equality in the gate is accepted throughout.
Since the comparison is pathwise, arbitrary permitted attack randomization
does not change it.
\end{proof}

\begin{proposition}[Finite suffix and half-gate upper bounds]
\label{prop:finite-upper}
Let $M\ge64$, $2\le r\le M+2$, and set
\[
 m=\min\{r-2,\lfloor M/2\rfloor\},\quad
 N=M-m,\quad L=1+\log\frac M{m+1}.
\]
From the histogram obtain $J'_u,C'_u$, and from the last $m+2$ symbols
obtain $T'_u=\sum_{i=M-m}^{M-1}\mathbf1\{Y'_i=u\}$.
Use the gate in Lemma~\ref{lem:gate} with $D'_u=O'_u=J'_u-T'_u$.
This $k,\beta$-blind estimator has all-start risk
\begin{equation}\label{eq:p-suffix-upper}
 R\le\min\left\{1,\frac{65536}{\sqrt M}
       +\frac{65536}{M\sqrt\beta}
       +(8448L+147744)\frac kM\right\}.
\end{equation}
Alternatively, use $D'_u=J'_{1,u}$, the count over the first
$\lfloor M/2\rfloor$ departures.  Its all-start risk is
\begin{equation}\label{eq:p-half-upper}
 R\le\min\left\{1,\frac{65536}{\sqrt M}
       +\frac{65536}{M\sqrt\beta}+2^{23}\frac kM\right\}.
\end{equation}
The latter estimator factors exactly through the histogram, protected
pair, and the one bit $\mathbf1\{512J'_{1,Y_M}\ge J'_{Y_M}\}$.
\end{proposition}
\begin{proof}
The last $m+2$ raw symbols reveal exactly $m$ departure pairs and the
terminal pair.  Thus $O'_u$ is precisely the count over indices
$0,\ldots,N-1$, even though their raw order is not observed.
In particular $r=2$ means $m=0$ and $O'=J'$; counting $r-1$ suffix
departures would incorrectly include the terminal pair.
Since $N\ge\lceil M/2\rceil$, the clean $O$ is at least $J_1$.
Lemma~\ref{lem:imbalance} therefore bounds the bias in
\eqref{eq:p-gate-reduction} by $65536/(M\sqrt\beta)$.
Lemma~\ref{lem:old-inverse}, with $d=4,c=1/1024$, bounds its reciprocal
term by $48k\cdot176L/M=8448Lk/M$.
Lemma~\ref{lem:occupation} bounds the two tails together by
\[
 \frac{16(6156+3078)k+32}{M+1}
 =\frac{147744k+32}{M+1}.
\]
Lemma~\ref{lem:empirical} and the remaining $32/(M+1)$ are bounded by
$65536/\sqrt M$.  The unit loss cap proves \eqref{eq:p-suffix-upper}.

For the half gate we require a different reciprocal estimate, valid for
any four-state chain and start:
\begin{equation}\label{eq:p-half-inverse}
 \E\frac{\mathbf1\{J_1\ge J/1024\}}{J+1}\le\frac{130944}{M}.
\end{equation}
To prove it, fix $u$ and condition on the entire first-half history,
whose terminal state is $v$ and whose count of $u$ is $j$.
Let $m_2=M-\lfloor M/2\rfloor$ be the remaining length.  Balance implies
$J_{2,u}\le1023j$.  In the suffix, the fixed-terminal event
$Y_M=u,J_{2,u}\le1023j$ is a subset of the actual-terminal low-occupation
event.  Lemma~\ref{lem:occupation} gives conditional probability at most
$16(1023j+1)/(m_2+1)$.  Since $(J_u+1)^{-1}\le(j+1)^{-1}$,
the conditional contribution is at most $16368/(m_2+1)$, including $j=0$.
Sum over four targets and use $m_2+1\ge M/2$ to prove
\eqref{eq:p-half-inverse}.  This argument uses a deterministic prefix,
not endpoint-conditioned independent samples; Lemma~\ref{lem:occupation}
is valid at every suffix length.
Apply Lemmas~\ref{lem:gate}, \ref{lem:imbalance}, and
\ref{lem:empirical} again.  The total corruption coefficient is
\[
 48\cdot130944+16(6156+3078)=6433056<2^{23},
\]
proving \eqref{eq:p-half-upper}.
The stated bit is exactly the gate decision on every corrupted word.
If $J'=0$ its value is one, but the decoder's separate $J'>0$ test still
returns zero.  Thus neither this boundary case nor threshold equality
requires another message.  All edits occur before the exact message is
formed; the bit channel itself is not being corrupted.

At $\beta=1/M$, \eqref{eq:p-suffix-upper} gives directly, for
$g_M(k)=\max\{M^{-1/2},k/M\}$,
\[
 R\le(131072+156192L)g_M(k)
 \le287264L g_M(k)<2^{19}L g_M(k).
\]
For $r=\max\{2,\lceil\alpha M\rceil\}$ with $0<\alpha\le1$,
$m+1=\min\{r-1,\lfloor M/2\rfloor+1\}\ge\alpha M/2$.
Indeed $r-1\ge\alpha M/2$ follows by separating $\alpha M\le2$
from $\alpha M>2$.  Hence $L\le1+\log(2/\alpha)$.
Likewise \eqref{eq:p-half-upper} becomes
$\min\{1,2^{17}/\sqrt M+2^{23}k/M\}$ at the critical certificate.
\end{proof}

\subsection{A count-only oracle when the budget is supplied}

\begin{proposition}[Critical known-budget upper]\label{prop:count-oracle}
At supplied certificate $\beta=1/M$, a predictor using only the exact
triple histogram, the protected pair, and the public budget $k$ satisfies,
for every source and initial law,
\[
 R\le191488\max\{M^{-1/2},k/M\},\qquad M\ge64.
\]
For $k\ge1$ its designated probability is
\[
 \widehat q_k=\frac{(C'-6k)_+}{J'+6k+2};
\]
for $k=0$ use $(C'+1)/(J'+2)$.  This budget-dependent estimator is used
only for the oracle comparison, not for the unknown-budget results.
\end{proposition}
\begin{proof}
\emph{A weighted occupation estimate.}
For an irreducible $d$-state chain with target masses $p_u$, let
$h_u=\max\{H_u,1/p_u\}$, $0\le q_u\le1$, and $W_u=p_uq_uh_u$.
For every $A>0$ and every initial law we first prove
\begin{equation}\label{eq:p-oracle-weighted-general}
 \E\min\left\{q_{Y_M},\frac A{J+2}\right\}
 \le384d^2\frac AM+\frac{192d}{M}\sum_uW_u.
\end{equation}
Fix a target $u$.  If $h_u>M$, the endpoint bound gives
$q_u\Pr(Y_M=u)\le4dp_uq_u(1+h_u/M)\le8dW_u/M$.
If $h_u\le M$, on $J_u\ge Mp_u/48$ the minimum is at most
$48A/(Mp_u)$, and $\Pr(Y_M=u)\le8dp_u$; this costs at most
$384dA/M$.  For the complementary event, put $n=\lfloor M/2\rfloor$
and $m_2=M-n$.  Since $n\ge M/3$,
$J_u<Mp_u/48$ implies $J_{1,u}<np_u/16$.
Conditioning on the complete prefix through $Y_n$, use the uniform
endpoint bound on the remaining $m_2$ transitions, and then the
unconditional lower tail on the prefix.  The weighted contribution is
at most
\[
 8dp_uq_u(1+2h_u/M)e^{-M/(24h_u)}
 =\frac{8dW_u}{M}(x+2)e^{-x/24}
 \le\frac{192dW_u}{M},\qquad x=M/h_u.
\]
The last step uses
$\sup_{x\ge0}(x+2)e^{-x/24}=24e^{-11/12}<24$.
Summing these two cases over targets proves
\eqref{eq:p-oracle-weighted-general}.
For full-support binary sources, Lemma~\ref{lem:hitting-scale} yields
\begin{equation}\label{eq:p-oracle-weighted-binary}
 \E\min\left\{q_{Y_M},\frac A{J+2}\right\}
 \le6144\frac AM+\frac{30720}{M\sqrt\beta}.
\end{equation}

For singleton support $00$, the $00$ contribution vanishes.  Use the
reflected chain from the proof of Lemma~\ref{lem:imbalance}.  On paths
ending active, original path weights are bounded above by reflected
path weights, with exactly the same terminal state and departure counts.
Thus the nonnegative minimum functional for targets $A,C$, with weights
$(b,0,\tau)$, is dominated pathwise in the same sum over paths.
Equation~\eqref{eq:p-reflected-weights} and
\eqref{eq:p-oracle-weighted-general}, with $d=3$, bound this part by
$3456A/M+1152/M$.  The original target $B$ contributes at most
$e\Pr(Y_M=B)\le4/(M+2)$ by \eqref{eq:p-next-absorption}.
An initial $00$ contributes zero.  For boundary $b$ or $\tau$, approximate
only the reflected chain by positive parameters: there are finitely many
length-$M$ paths, and each summand is continuous in probabilities and
weights, including its minimum.  The uniform bound therefore passes to
the limit.  Symbol exchange handles singleton $11$.
Since $\beta\le1$, the larger bound
\eqref{eq:p-oracle-weighted-binary} covers both singleton cases and
arbitrary initial mixtures as well.

\emph{A pathwise lower envelope.}
For $k\ge1$ put $B_k=6k$, $D=J+2$, $D'=J'+B_k+2$, and
$c'=(C'-B_k)_+$.  The protected-edit perturbation bounds give
\[
 D'\ge D,\quad0\le c'\le C,\quad
 0\le C-c'\le2B_k,\quad0\le D'-D\le2B_k.
\]
Thus $0\le\widehat q_k\le C/(J+2)$, and
\[
 \frac CD-\frac{c'}{D'}
 =\frac{C-c'}D+\frac{c'}{D'}\frac{D'-D}D
 \le\frac{4B_k}D.
\]
For $J>0$, putting $p=C/J$ gives
\[
 0\le p-\widehat q_k
 \le\min\left\{p,\frac{4B_k+2}{J+2}\right\}.
\]
Since $x\mapsto\min\{x,c\}$ is $1$-Lipschitz on nonnegative reals,
with $A=24k+2$ the triangle inequality implies
\begin{equation}\label{eq:p-oracle-pathwise}
 |\widehat q_k-q|
 \le2\mathbf1\{J>0\}|C/J-q|
       +\min\left\{q,\frac A{J+2}\right\}.
\end{equation}
When $J=0$, $C=0$ and $\widehat q_k=0$; since $A\ge2$, the last
term equals $q$, so the inequality still holds without defining $C/J$.
This is valid for every permitted edit realization.
Using Lemma~\ref{lem:empirical},
\eqref{eq:p-oracle-weighted-binary}, $\beta=1/M$, and $A\le26k$ gives
\[
 R\le\frac{31744}{\sqrt M}+159744\frac kM
 \le191488\max\{M^{-1/2},k/M\},\qquad k\ge1.
\]

For $k=0$ the word is unedited.  On $J>0$, the add-one probability
is within $1/(J+2)$ of $C/J$; on $J=0$ pay unit loss.
Lemma~\ref{lem:occupation} gives $\Pr(J=0)\le16/(M+1)$.
For completeness, integer-threshold layer cake gives
\[
 \E\frac1{J+2}
 \le\int_0^1\min\left\{1,\frac{16}{(M+1)t}\right\}\,dt
 \le\frac8{\sqrt{M+1}}.
\]
Indeed a nonempty event $J<t^{-1}-2$ uses threshold
$\lceil t^{-1}-2\rceil\le t^{-1}$; for the last inequality use
$\min\{1,z\}\le\sqrt z$ and integrate $t^{-1/2}$.
Together with Lemma~\ref{lem:empirical} this is at most
$536/\sqrt M$, below the asserted constant.  No positional message or
raw-window input is used in either budget branch.
\end{proof}

\section{Observation separation and lower bounds}
\label{sec:negative-proofs}

Throughout this section a word has symbols $X_0,\ldots,X_{M+1}$,
pair states $Y_i=(X_i,X_{i+1})$, and departures $0\le i<M$.
Write $J_{ab}=\sum_{i<M}\ind\{Y_i=ab\}$,
$J_{\rm cross}=J_{01}+J_{10}$, and $N_{abc}=H_{abc}(X)$.
Thus $Y_M$ is a terminal pair, not a departure. All risks below include
independent predictor randomness. The observation $\mathcal T_r$ consists
of the eight triple counts and the actual first and last $r$ symbols.

\subsection{An exact occupation mixture}

For $0<a,t<1$, let $Q_a$ be the stationary law with pair transition matrix,
in the order $00,01,10,11$,
\[
 P_a=\begin{pmatrix}
 1-a&a&0&0\\
 0&0&1/2&1/2\\
 1/2&1/2&0&0\\
 0&0&t&1-t
 \end{pmatrix},\qquad
 \Pi_a=\frac{(t/a,2t,2t,1)}{1+t/a+4t}.
\]
In particular the stationary initial factor is
$\Pi_a(11)=a/[t+(1+4t)a]$.

\begin{lemma}[Occupation-mixture domination]
\label{lem:mixture}
Put $F=\{Y_0=11,Y_M=00\}$ and $U=N_{001}$. For every integer $j$,
every integer $u_0\ge1$, every interval $I\subset(0,1)$, and every
additional event $G$ (possibly depending on $a$),
\[
 \int_I Q_a(F,J_{00}=j,U\ge u_0,G)\,\frac{da}{a}
 \le \frac{t}{u_0}.
\]
\end{lemma}
\begin{proof}
Delete every transition departing from $00$, both loops and exits.
For $J_{00}=j$, exactly $L=M-j$ transitions remain. Their last transition
is $10\to00$. Before that transition the retained path is on
$B=01,C=10,D=11$. A retained $C\to00$ followed by an erased exit
$00\to B$ is represented by a marked $C\to B$ of weight $1/2$;
a direct $C\to B$ is unmarked and also has weight $1/2$. The erased
exit's factor $a$ is restored separately. All other retained transitions
keep their probabilities. The marked transition matrix is therefore
\[
 K(z)=\begin{pmatrix}
 0&1/2&1/2\\ (1+z)/2&0&0\\ 0&t&1-t
 \end{pmatrix}.
\]
For $L\ge1$, define
\[
 h_{L,u}=\tfrac12[z^u](K(z)^{L-1})_{D,C},
\]
and set $h_{L,u}=0$ for $L\le0$ or $u<0$. The final transition
$C\to00$ has weight $1/2$ and is not marked. A skeleton with $u$ marks
has $u$ completed visits to $00$ and one terminal visit. Distributing
its $j-u$ loops among these $u+1$ holding slots gives exactly
$\binom ju$ weak compositions. Each composition and marked skeleton
recover one word, and deletion recovers both uniquely. Consequently,
for $0\le u\le j$,
\begin{equation}
 Q_a(F,J_{00}=j,U=u)
 =\Pi_a(11)\binom ju a^u(1-a)^{j-u}h_{M-j,u}.
 \label{eq:marked-word-identity}
\end{equation}
Infeasible cases have zero mass. The terminal holding slot is present
even though $Y_M$ is not a departure; there is no initial $00$ slot.

For $u\ge1$, the beta integral gives
$\binom ju\mathrm B(u,j-u+1)=1/u$. Thus, if
$A_{j,u}$ has the $\operatorname{Beta}(u,j-u+1)$ law,
\begin{align}
 \int_0^1 Q_a(F,J_{00}=j,U=u)\,\frac{da}{a}
 &=\frac{h_{M-j,u}}{u}\E\Pi_{A_{j,u}}(11)
 \le\frac{h_{M-j,u}}{u}. \label{eq:stationary-beta-factor}
\end{align}
The expectation is part of the equality: stationarity has not been
cancelled. For $u=0$ the corresponding factor is instead
$h_{M-j,0}\int_0^1(1-a)^j/[t+(1+4t)a]\,da$, and no $1/u$ formula
is used. This is why $u_0>0$ is required.

To sum over marks at the exact outside horizon, put $R=K(1)$.
This matrix is stochastic. Started at $D$, let $b_n,c_n,d_n$ be its
probabilities of $B,C,D$. Then
\[
 c_0=0,\quad c_1=t,\quad b_n=c_{n-1}\ (n\ge1),\qquad
 c_n=\tfrac12 c_{n-2}+td_{n-1}\ (n\ge2).
\]
Induction gives $c_n\le2t$. Hence, for each $L\ge1$,
\[
 \sum_{u\ge0}h_{L,u}=\tfrac12 c_{L-1}\le t.
\]
First drop $G$ and enlarge $I$ to $(0,1)$ in the nonnegative integral.
Then sum \eqref{eq:stationary-beta-factor} over $u\ge u_0$ to obtain
\[
 \int_I Q_a(F,J_{00}=j,U\ge u_0,G)\,\frac{da}{a}
 \le\sum_{u\ge u_0}\frac{h_{M-j,u}}u
 \le\frac1{u_0}\sum_{u\ge0}h_{M-j,u}\le\frac t{u_0}.
\]
No conditioning on a prescribed number of outside marks or replacement
of the horizon $M-j$ is involved.
\end{proof}

\subsection{Relocating loops while preserving sampled windows}

\begin{lemma}[Interior loop relocation]
\label{lem:relocation}
Suppose $Y_0=11,Y_M=00$, $2\le r$, and $2r\le M+2$.
Put $l=r-2$, $h=M-r+2$. If the interior vertices
$Y_l,\ldots,Y_h$ contain a visit to $11$ strictly before a visit to $00$,
there is a deterministic transformation $C_r$ preserving the complete
triple histogram and the actual first and last $r$ symbols, such that,
with $j=J_{00}$,
\begin{equation}
 d_H\bigl(C_r(X),1^{M+2-j}0^j\bigr)
 \le3J_{\rm cross}+4r-5.
 \label{eq:relocation-distance}
\end{equation}
If at least one of the two movable diagonal-loop types is absent, the
same conclusion holds without the ordered-visit condition.
\end{lemma}
\begin{proof}
Freeze departures $i<l$ and $i\ge h$. These are $2r-4$ edges and
spell exactly the sampled boundary symbols. From the interior path
remove only the $111$ and $000$ loops, in numbers $A_1$ and $A_0$,
retaining every other edge in order. Let $S_0,\ldots,S_L$ be its vertices.
Deleting self-loops neither removes a distinct-state visit nor reverses
its order. Choose the first retained $11$ and the first subsequent
retained $00$; insert all $A_1$ loops at the former and all $A_0$ at
the latter. The choice is deterministic. Deleting or inserting a loop
preserves its attachment pair. Reinserting exactly the removed number
of edges preserves the absolute suffix attachment index $h$, so the
frozen edge labels and all sampled boundary symbols are unchanged.
Every triple label is retained or reinserted once, proving histogram
preservation. Unfrozen boundary-crossing triples need not keep their
individual positions.

Here is the required positional count argument. A binary word of length
$N$ with a distinguished block of $A_1$ ones before a distinguished
block of $A_0$ zeros has $o$ other ones and $z$ other zeros. If
$j=A_0+q$, $q\ge0$, then
\begin{equation}
 d_H(\text{word},1^{N-j}0^j)
 \le\max\{z+q,\,2o+z-q\}. \label{eq:two-block-distance}
\end{equation}
Indeed, write $z_-$ for zeros before the target cut. The exact mismatch
count is $2z_-+j-(A_0+z)$. If a distinguished one lies after the cut,
all distinguished zeros lie after the cut and $z_-\le z$, giving
$z+q$. Otherwise every one after the cut is among the $o$ other ones;
writing $o_+$ for their number, the other exact identity is
$2o_++(A_0+z)-j$, at most $2o+z-q$. This proof does not presume any
order among the remaining symbols.

Put $u=N_{001}$, $v=N_{011}$, $p=N_{010}$, $w=N_{101}$, and
$d=J_{\rm cross}$. Endpoint flow gives
\[
 N_{100}=u+1,\qquad N_{110}=v+1,\qquad
 d=p+v+u+1+w.
\]
Let $f_1,f_0$ be the frozen $111,000$ loop counts, so
$f_1+f_0\le2r-4$. The initial pair contributes two other ones.
Counting the appended symbol of every remaining edge gives exactly
\begin{align*}
 o&=2+u+v+w+f_1=d+1-p+f_1,\\
 z&=p+u+v+2+f_0=d+1-w+f_0,\\
 q&=j-A_0=u+f_0\ge0.
\end{align*}
Since $u\le d-1$ and $d\ge1$,
\begin{align*}
 z+q&=d+1-w+u+2f_0\le2d+2f_0,\\
 2o+z-q&=3d+3-2p-w-u+2f_1\le3d+3+2f_1.
\end{align*}
Each is at most $3d+4r-5$, proving \eqref{eq:relocation-distance}.
If only one movable loop type occurs, put its block at any retained
visit to its state, which exists because the loops were deleted there.
The same two-block argument treats the absent block as empty; if both
are absent use the identity transformation.

Finally, $j$ is the sampled word's actual departure count. The comparison
word has instead
\begin{equation}
 J_{00}(1^{M+2-j}0^j)=j-2\quad(2\le j\le M).
 \label{eq:null-suffix-shift}
\end{equation}
For this range its final pair and the transformed word's final pair are
both $00$, so all their Hamming differences are in the editable first
$M$ coordinates.
\end{proof}

\subsection{A finite matched-observation lower bound}

\begin{proposition}[Finite critical lower bound]
\label{prop:finite-lower}
Let
\[
 c=2^{-16},\quad A=2^{12},\quad t=\frac{ck}{M},\quad
 u_0=\lfloor ck/4\rfloor,\quad T=\lfloor M/8\rfloor,\quad
 L_-=\log\frac1{4At}.
\]
Assume integer $M,k,r$ satisfy
\begin{equation}
 \begin{gathered}
 M\ge64,\quad 2\le r\le k/128,\quad k\le M/8,\quad ck\ge1024,\\
 At<1/4,\qquad \frac M{k^2}\le\frac{Ac^2}{2^{20}}.
 \end{gathered}\label{eq:finite-lower-domain}
\end{equation}
For any fixed predictor based on $\mathcal T_r$, let $S$ be its supremum
clean risk over the stationary laws $Q_a$, $At\le a\le1/4$.
Every such source satisfies the all-state certificate with supplied
$\beta=1/M$. There is a stationary source with the same supplied
certificate and one source-independent strict-$k$ corruption kernel
for which
\begin{equation}
 R_{\rm null}\ge\frac{u_0}{8Mt}
       \left[\left(\frac t{16}-S\right)L_-\right]_+.
 \label{eq:finite-lower}
\end{equation}
The kernel protects the final pair on every input and every output.
\end{proposition}
\begin{proof}
We give the source, probability, and kernel calculations separately.

\paragraph{Source legality and all sixteen certificate entries.}
Throughout $t\le c/8$, $a\ge At\ge t$, and $a\le1/4$. Set
$D=1+t/a+4t$. Then
$1\le D\le1+1/A+c/2<4/3<2$.
The flow equalities
$\Pi_a(00)a=\Pi_a(10)/2$ and
$\Pi_a(11)t=\Pi_a(01)/2$, together with
$\Pi_a(01)=\Pi_a(10)$, verify stationarity directly. Multiplication gives
\begin{equation}
 P_a^2=\begin{pmatrix}
 (1-a)^2&(1-a)a&a/2&a/2\\
 1/4&1/4&t/2&(1-t)/2\\
 (1-a)/2&a/2&1/4&1/4\\
 t/2&t/2&(1-t)t&(1-t)^2
 \end{pmatrix}.
 \label{eq:all-sixteen}
\end{equation}
Every displayed entry is at least $t/2$: use $a\ge t$ and
$a,t\le1/4$ for the products and all other entries are immediate.
Since $t/2=ck/(2M)\ge1/M$, every starting and destination state obeys
$P_a^2(u,v)\ge M^{-1}\Pi_a(v)$. Also $a\ge At\ge2/M$, so
$a-\beta\ge a/2$.

\paragraph{Common-atom covariance and coupling bounds.}
The four paths to $11$ of length three, from the states in order, have
probabilities
\[
 (1-a)a/2,\qquad (1-t)^2/2,\qquad a/4,\qquad(1-t)^3.
\]
These arise respectively from paths $00,00,01,11$;
$01,11,11,11$; $10,00,01,11$; and $11,11,11,11$.
Each is at least $\lambda=a/4$, so
$P_a^3(x,\cdot)\ge\lambda\delta_{11}$ for every $x$.
Split each three-step transition into this common atom and a residual
probability law. Shared independent Bernoulli trials of probability
$\lambda$ couple any two chains at the first common atom. Fill in each
three-step block by its conditional bridge; this preserves the full
one-step path law. After meeting, use identical transitions. Thus the
contraction coefficient at lag $h$ is bounded by
\[
 \alpha_h=(1-\lambda)^{\lfloor h/3\rfloor}.
\]

For a stationary chain $W$, a departure-state-set count with stationary
mass $p$, or a count of a specified directed edge with stationary
probability $p$, satisfies
\begin{equation}
 \operatorname{Var}(V)\le\frac{7Mp}{\lambda}.
 \label{eq:common-atom-variance}
\end{equation}
For an edge indicator $I_i$ and $h\ge1$, condition on the state after
edge $i$. The conditional probability of edge $i+h$ is a $[0,1]$-valued
function after $h-1$ transitions, whose deviation from its stationary
value is at most $\alpha_{h-1}$. Therefore
$|\operatorname{Cov}(I_i,I_{i+h})|\le p\alpha_{h-1}$.
For a state-set indicator use $\alpha_h\le\alpha_{h-1}$ instead.
Since $\sum_{h\ge1}\alpha_{h-1}=3/\lambda$, expansion of the variance
gives $Mp(1+6/\lambda)\le7Mp/\lambda$, as claimed.

Couple a chain $Z$ started at $00$ to stationary $W$ by the same block
construction. Its meeting boundary $\sigma$ has
\[
 \E\sigma\le3/\lambda=12/a.
\]
Every departure-state or edge count in the two paths differs by at most
$\sigma$ pathwise: after boundary $\sigma$ the paths agree, and there
are at most $\sigma$ earlier departures. No independence of $\sigma$
and the stationary counts is needed.

\paragraph{Stationary reversal and count corrections.}
Let $\mathcal R(ab)=ba$. The stationary triple identity
\[
 \Pi_a(ab)P_a(ab,bc)=\Pi_a(cb)P_a(cb,ba)
\]
holds for every triple: the only nonpalindromic cases are $001,100$
and $011,110$, given by the two flow identities above. It implies
$\mathcal R P_a^*\mathcal R=P_a$, where $P_a^*$ is the stationary
reverse pair kernel. Reversing the symbols of a stationary word thus
preserves its law. In particular, with
$F=\{Y_0=11,Y_M=00\}$ and $E=\{Z_M=11\}$,
\[
 Q_a(F)=\Pi_a(00)P_{00}(E),
\]
and conditional on $F$, the reversed pair path
$Z_i=\mathcal RY_{M-i}$ has law $P_{00}(\cdot\mid E)$.
For $M\ge64$, $\lfloor M/3\rfloor\ge M/4$, and therefore
\[
 \alpha_M\le e^{-aM/16}\le\frac{16}{aM}
 \le\frac{16}{Ack}\le\frac14.
\]
Since $\Pi_a(11)=1/D\ge3/4$, we have $P_{00}(E)\ge1/2$.

Put $V(W)=N_{001}(W)$ and $x=ck$. Its stationary mean is
$Mt/D\in[x/2,x]$. Consider the four conditions
\begin{equation}
 \sigma\le x/16,\quad |V(W)-Mt/D|\le x/8,\quad
 J_{\rm cross}(W)\le k/32,\quad J_{00}(W)\le M/16.
 \label{eq:preliminary-good}
\end{equation}
Their failure probabilities, respectively, are at most
\begin{equation}
 \frac{192M}{Ac^2k^2},\quad
 \frac{1792M}{Ac^2k^2},\quad
 \frac{458752M}{Ak^2},\quad \frac{16}{A}.
 \label{eq:four-failures}
\end{equation}
The first follows from $\E\sigma\le12/a$ and $a\ge Ack/M$.
For the second, \eqref{eq:common-atom-variance} with $p=t/D$ gives
$\operatorname{Var}V\le28Mt/a$; divide by $(ck/8)^2$.
For the third, the mean is $4Mt/D\le4ck\le k/64$, and
$\operatorname{Var}J_{\rm cross}\le112Mt/a\le112M/A$;
divide by $(k/64)^2$. For the last use
$\Pi_a(00)=t/(aD)\le1/A$ and Markov's inequality.
Conditioning the joint coupling on $E$ costs at most two. Under
\eqref{eq:finite-lower-domain}, twice the sum in
\eqref{eq:four-failures} is bounded by
\[
 2\left(\frac{1984}{2^{20}}
       +\frac{458752c^2}{2^{20}}+\frac{16}{A}\right)<\frac1{16}.
\]

On $F$, the exact reversal identities are
\begin{equation}
 \begin{split}
 J_{00}(Z)&=J_{00}(Y)+1,\\
 J_{\rm cross}(Z)&=J_{\rm cross}(Y),\\
 N_{001}(Z)&=N_{100}(Y)=N_{001}(Y)+1.
 \end{split}\label{eq:reverse-count-shifts}
\end{equation}
The first compares arrival indices $1,\ldots,M$ with departure indices
$0,\ldots,M-1$; the second loses no mixed endpoint; the third combines
triple reversal with entry-minus-exit flow at $00$.
Thus \eqref{eq:preliminary-good}, the pathwise coupling bound, and
$x\ge1024$ imply for the original path
\begin{align*}
 U&\ge x/2-x/8-x/16-1\ge x/4\ge u_0,\\
 J_{\rm cross}&\le k/32+ck/16\le k/16,\\
 2\le u_0&\le J_{00}\le M/16+ck/16\le M/8.
\end{align*}
The integer upper count bound is $J_{00}\le T$.

\paragraph{Ordered visits in the interior.}
Let $l=r-2$ and $h=M-r+2$. Since $r\le M/1024$,
$h-l=M-2r+4\ge M/2$. Write $B$ for a visit to $11$ before a later
visit to $00$ within $Y_l,\ldots,Y_h$. Because $M-h=l$ and $M-l=h$,
reversal sends this event precisely to a visit to $00$ before a later
visit to $11$ within $Z_l,\ldots,Z_h$.

The ordinary expected hitting times of $00$ and $11$, zero at their
targets and listed in state order, are
\begin{align*}
 h_{00}&=(0,\,4+1/t,\,3+1/(2t),\,3+3/(2t)),\\
 h_{11}&=(3+3/(2a),\,3+1/(2a),\,4+1/a,\,0).
\end{align*}
For completeness, the first vector follows by solving
$h_{01}=1+(h_{10}+h_{11})/2$,
$h_{10}=1+h_{01}/2$, and
$h_{11}=1+th_{10}+(1-t)h_{11}$, with $h_{00}=0$.
For the second solve
$h_{00}=1+(1-a)h_{00}+ah_{01}$,
$h_{01}=1+h_{10}/2$, and
$h_{10}=1+(h_{00}+h_{01})/2$, with $h_{11}=0$.
On $a,t\le1/4$ the two maxima are
$3+3/(2t)$ and $3+3/(2a)$. From any state at time $l$, hitting $00$
and then $11$ therefore has expected additional time at most
\[
 6+\frac{3}{2t}+\frac{3}{2a}\le\frac6t.
\]
This uses the strong Markov property at the first hit: partition by its
finite time and state, apply the usual Markov property on each event,
and sum. Distinct targets make the ordering strict. Markov's inequality
over $h-l\ge M/2$ bounds the unconditioned failure probability by
$12/(tM)$. Dividing by $P_{00}(E)\ge1/2$ gives
\[
 Q_a(B^c\mid F)\le\frac{24}{ck}\le\frac3{128}.
\]
This estimate is uniform in the time-$l$ state; it assumes no
independence from the final state or the counts.

Let $G$ be the intersection of $B$ with
$U\ge u_0$, $J_{\rm cross}\le k/16$, and $2\le J_{00}\le T$.
The preceding count and hitting bounds give
$Q_a(G\mid F)\ge1-1/16-3/128>1/2$. Retaining the rare endpoint mass,
\begin{equation}
 Q_a(F\cap G)\ge\Pi_a(00)/4,\qquad
 (a-\beta)Q_a(F\cap G)\ge t/16.
 \label{eq:good-event-mass}
\end{equation}
Indeed $P_{00}(E)\ge1/2$, $a-\beta\ge a/2$, and
$a\Pi_a(00)=t/D\ge t/2$.

\paragraph{The null and its full word masses.}
Take a stationary first-order binary source with
$T_*(1,0)=q_*=1/(2M)$ and $T_*(0,1)=\beta=1/M$.
Its symbol masses are $\pi_*(0)=1/3$, $\pi_*(1)=2/3$.
Its stationary pair law is $\Pi_*(s,z)=\pi_*(s)T_*(s,z)$, and
for all starting and destination pair states,
\[
 \frac{P_*^2((u,v),(s,z))}{\Pi_*(s,z)}
 =\frac{T_*(v,s)}{\pi_*(s)}.
\]
For $(v,s)=(0,0),(0,1),(1,0),(1,1)$ these four ratios are
$3(1-\beta),\,3\beta/2,\,3q_*,\,3(1-q_*)/2$, respectively.
They are all at least $\beta+q_*\ge\beta$, checking every pair-state
certificate entry. The target at terminal $00$ is Bernoulli$(\beta)$,
not Bernoulli$(q_*)$.
For $2\le j\le T$ put $x^{(j)}=1^{M+2-j}0^j$.
These disjoint words have full stationary masses
\begin{equation}
 v_j=\frac23(1-q_*)^{M+1-j}q_*(1-\beta)^{j-1}
 \ge\frac1{8M}. \label{eq:null-word-masses}
\end{equation}
Bernoulli's inequality gives the prefix factor at least $1/2$ and
the suffix factor at least $7/8$, so in fact the product is at least
$7/(48M)>1/(8M)$. These masses include the initial symbol and every
one of the $M+1$ transitions.

\paragraph{A predeclared strict kernel with identity rejection.}
For $2\le j\le T$, define a finite measure on words by
\[
 \nu_j(dy)=\int_{At}^{1/4}
       Q_a(dy;F,U\ge u_0,J_{00}=j)\,\frac{da}{a},
 \qquad \mu_j=\nu_j(\text{all words}).
\]
Lemma~\ref{lem:mixture} gives $\mu_j\le t/u_0$.
For $\mu_j>0$, normalization $\nu_j/\mu_j$ is a probability law.
These laws depend only on declared $M,k,r$ and the auxiliary family,
not the true source, its parameters, its initial law, or a predictor seed.

On input $x^{(j)}$ with $\mu_j>0$, draw $Y$ once from
$\nu_j/\mu_j$. If $B$ holds, form the deterministic $C_r(Y)$ of
Lemma~\ref{lem:relocation}. Accept it only if its actual Hamming
distance from $x^{(j)}$ is at most $k$ and its final pair equals that
of $x^{(j)}$. Otherwise output the unchanged input. On zero-mass labels
and all other input words, also output the unchanged input. This defines
one source-independent kernel on all words, with a strict pointwise
budget and protected final pair, not merely an average budget.
It uses one draw, not repeated sampling until acceptance.

On $G$, Lemma~\ref{lem:relocation} bounds the distance by
$3k/16+4r\le3k/16+k/32<k$, so every check succeeds, and
\[
 \mathcal T_r(C_r(Y))=\mathcal T_r(Y).
\]
The same statistic identity holds on every transformed acceptance;
it is not asserted on identity rejections. These are the sampled
alternative's boundary windows, not necessarily the null input's
windows. Input-window edits are allowed and included in the Hamming
check. The suffix label obeys exactly
$J_{00}(x^{(j)})=j-2$ and $J_{00}(Y)=j$.

\paragraph{Nonnegative-loss ordering.}
For a sampled word $y$, let $\ell_a(y)$ be the predictor's expected
absolute error against $a$ on $\mathcal T_r(y)$, and let $\ell_*(y)$
be its corresponding error against $\beta$. On $F$ the clean actual
terminal target is $a$; the triangle inequality, also after averaging
independent decoder randomness, gives
$\ell_*(y)\ge a-\beta-\ell_a(y)$.
For $\mu_j>0$, \eqref{eq:null-word-masses} gives
\[
 \frac{v_j}{\mu_j}\ge\rho:=\frac{u_0}{8Mt}.
\]
First discard nonnegative losses from other null inputs and rejected
proposals. On the remaining nonnegative accepted losses, replace
$v_j/\mu_j$ by $\rho$. Next restrict to $G$, where acceptance and
statistic matching hold. Only then use the signed triangle bound:
\begin{align*}
 R_{\rm null}
 &\ge\rho\int_{At}^{1/4}
          \E_{Q_a}[\ind_{F\cap G}\ell_*(Y)]\,\frac{da}{a}\\
 &\ge\rho\int_{At}^{1/4}
    \left((a-\beta)Q_a(F\cap G)
          -\E_{Q_a}[\ind_{F\cap G}\ell_a(Y)]\right)\frac{da}{a}\\
 &\ge\rho\left(\frac t{16}-S\right)L_-.
\end{align*}
The restricted alternative loss is bounded by its unconditional clean
risk, hence by $S$. Risk is nonnegative, giving \eqref{eq:finite-lower}.
In particular, no lower mass multiplier has been placed in front of a
possibly negative quantity, and no endpoint-conditioned loss has been
substituted for the unconditional target risk.
\end{proof}

\subsection{Consequences for widths and supplemental messages}

\begin{proposition}[Width and fixed-fraction consequences]
\label{prop:lower-consequences}
For a prescribed sequence $2\le r_M\le M+2$, a finite uniform factor
relative to $g_M(k)=\max\{M^{-1/2},k/M\}$ exists if and only if
$\liminf_{M\to\infty}r_M/M>0$, for stationary or arbitrary starts.
For each fixed $0<\alpha\le1/4096$ and
$r_M=\max\{2,\lceil\alpha M\rceil\}$,
\[
 2^{-26}\log\frac1{32\alpha}
 \le C_{\rm stat}(\alpha)
 \le C_{\rm all}(\alpha)
 \le2^{19}\left(1+\log\frac2\alpha\right).
\]
With exact counts and protected terminal pair, optionally also the
initial pair, a post-corruption full-word encoder needs at least two
supplemental symbols to attain a finite uniform factor; two suffice
when its message is exact and fault-free.
\end{proposition}
\begin{proof}
Suppose first that the liminf is zero. Select an increasing subsequence
on which $r_M/M\to0$, and on that subsequence put
\[
 b_M=\max\{r_M,\sqrt M\},\qquad
 k_M=\lfloor\sqrt{Mb_M}\rfloor.
\]
Then $b_M=o(M)$, $k_M/M\to0$, $k_M/b_M\to\infty$,
$r_M/k_M\to0$, and $k_M/\sqrt M\to\infty$. Moreover
$M/k_M^2\to0$, so every condition in
\eqref{eq:finite-lower-domain} holds eventually. Here
$t_M=ck_M/M$, $\rho=u_0/(8Mt_M)$, $L_-\to\infty$, and
$\rho\to1/32$.
If one $k$-blind family had a finite admissible factor $C$, its clean
risk on every alternative at this same $M$ and supplied certificate
would satisfy $S_M\le C/\sqrt M=o(t_M)$. For example, eventually
$\rho\ge1/64$ and $t_M/16-S_M\ge t_M/32$, so
\[
 \frac{R_{\rm null}}{g_M(k_M)}
 \ge\frac{c}{2048}L_-\longrightarrow\infty,
\]
since $g_M(k_M)=k_M/M$ eventually. This contradicts $C$.
The alternatives may depend on the tested budget, but the predictor
and its clean guarantee are the same map at the same horizon. These
stationary witnesses also rule out an all-start guarantee.

Conversely, if the liminf is positive, choose $\delta>0$ and $M_0$
so $r_M\ge\delta M$ for $M\ge M_0$. Put
$m_M=\min\{r_M-2,\lfloor M/2\rfloor\}$.
Then $(m_M+1)/M$ is bounded away from zero for all sufficiently large
$M$. The suffix estimate of Proposition~\ref{prop:finite-upper} is at
most $2^{19}[1+\log(M/(m_M+1))]g_M(k)$ at $\beta=1/M$, so gives
one uniform all-start factor eventually. A finite initial set of
horizons is harmless: any prediction has risk at most one and
$g_M(k)\ge M^{-1/2}$. This proves the exact criterion without
monotonicity of the widths.

Now fix $\alpha>0$ in the stated interval before sending $M$ to
infinity. For any family with finite admissible stationary factor $C$,
put $k=128r_M$ at sufficiently large horizons. Then
\[
 \frac kM\to128\alpha\le\frac1{32},\quad
 t\to128c\alpha>0,\quad At\to8\alpha<\frac14,
\]
all finite-domain conditions eventually hold, and
\[
 L_-\to\log\frac1{32\alpha},\qquad
 \frac{u_0}{8Mt}\to\frac1{32},\qquad
 S/t\le\frac{C}{t\sqrt M}\to0.
\]
For these budgets $k\ge\sqrt M$ eventually. Using the deliberately
loose $g_M(k)\le2k/M$ in \eqref{eq:finite-lower}, the positive limiting
bracket gives
\[
 C\ge\liminf_{M\to\infty}\frac{R_{\rm null}}{g_M(k)}
 \ge\frac1{32}\frac{c}{32}\log\frac1{32\alpha}
 =2^{-26}\log\frac1{32\alpha}.
\]
This holds for every admissible finite factor, so taking the infimum
does not require an attained optimum. The upper follows from
Proposition~\ref{prop:finite-upper}: $m_M+1\ge\alpha M/2$, including
$r_M=2$, hence $1+\log(M/(m_M+1))\le1+\log(2/\alpha)$.
The stationary class is a subclass of the arbitrary-start class.
These arguments take the horizon limit for each fixed fraction and
factor; they neither exchange limits nor give an every-budget
logarithmic risk lower bound.

Finally, a one-symbol supplemental alphabet is constant, so leaves
only counts and the supplied endpoint pair or pairs. At $r_M=2$ even
both endpoint pairs are insufficient by the necessity just proved.
The exact one-bit factorization in
Proposition~\ref{prop:finite-upper} supplies two symbols and has
all-start risk at most $(2^{17}+2^{23})g_M(k)$. This concerns a
fault-free message computed by scanning the whole already corrupted
word, not one raw coordinate or a message subject to further corruption.
\end{proof}

\subsection{Elementary same-information oracle lower floors}

\begin{lemma}[Critical oracle lower floors]
\label{lem:oracle-floor}
For every $M\ge64$ and $0\le k\le M$, the stationary minimax risk
with supplied $\beta=1/M$ is at least $g_M(k)/96$, even if the predictor
knows $k$ and observes the entire corrupted word. Consequently this
lower bound holds for each restricted observation considered here and
for the arbitrary-start minimax risk.
\end{lemma}
\begin{proof}
We give separate clean and edit witnesses; taking their maximum is
permitted in a minimax lower bound.

\paragraph{Clean floor.}
Put $N=M+2$ and use stationary IID Bernoulli sources
\[
 p_\pm=\tfrac12\pm\frac1{16\sqrt N}.
\]
Their pair chains satisfy $P^2(u,v)=\Pi(v)$ for every row, hence the
supplied certificate. The target is $p_\pm$ at every terminal pair.
The elementary Bernoulli divergence bound
\[
 d(p\Vert q)\le\frac{(p-q)^2}{q(1-q)}
\]
follows by applying $\log x\le x-1$ to its two summands.
With $\Delta=p_+-p_-=1/(8\sqrt N)$, it gives one-letter divergence
at most $1/(8N)$. Independence makes the full-word divergence at most
$1/8$, and their total variation distance is at most $1/4$.
For clarity, this last inequality follows without a statistical
independence assertion about a selected row: for finite probability
laws $P,Q$, group their masses on
$A=\{x:P(x)\ge Q(x)\}$ and its complement. The log-sum inequality
(Jensen's inequality for $x\log x$) gives
$D(P\Vert Q)\ge d(P(A)\Vert Q(A))$.
For fixed $q$, the second derivative in $p$ of $d(p\Vert q)$ is
$1/[p(1-p)]\ge4$; its value and first derivative vanish at $p=q$.
Thus $D(P\Vert Q)\ge2(P(A)-Q(A))^2=2\TV(P,Q)^2$.

Testing by the nearer target makes each error incur prediction loss
at least $\Delta/2$. Every possibly randomized test has average error
at least $(1-\TV(P,Q))/2$, since the difference of its acceptance
probabilities is at most total variation. Thus the average prediction
risk is at least
\[
 \frac\Delta4(1-\TV(P,Q))
 \ge\frac3{128\sqrt N}\ge\frac1{64\sqrt M}.
\]
Identity corruption is allowed at every $k$, including for a
known-budget predictor.

\paragraph{Edit floor.}
For $k\ge1$, use the uniform prior on IID Bernoulli parameters
$p\in[1/4,3/4]$. Put $b=k+1$. If the first $M$ symbols contain $n$
ones, retain a uniformly chosen set of
$r=b\lfloor n/b\rfloor$ of those ones, turning the others to zero.
Leave the final two symbols $U$ unchanged. This single kernel is
source-independent and changes $n-r\le b-1=k$ editable bits on every
outcome.
For a specified retained set $T$ of size $r$, conditional on $n$,
\[
 \Pr(T\text{ retained}\mid n)
 =\frac{\binom{M-r}{n-r}}
        {\binom Mn\binom nr}
 =\frac1{\binom Mr}.
\]
Hence conditional on $r,U$ the output arrangement is uniform and
independent of $p$, even after mixing the possible $n$ in a bucket.
The full observation factors through at most
\[
 D=4\bigl(\lfloor M/b\rfloor+1\bigr)
\]
labels $(r,U)$ followed by a $p$-independent channel. The four protected
pair possibilities are counted, not conditioned away.

Realize that channel using a seed independent of $p$ and of the label,
and likewise use an independent decoder seed. Freeze both seeds.
There are at most $D$ possible prediction values. For the uniform
prior on an interval of length $\ell=1/2$, the union of radius-$v$
neighborhoods of those values has prior mass at most $2Dv/\ell$,
regardless of how a label is randomly selected conditional on $p$.
Outside that union every possible prediction has error greater than
$v$. Integrating this probability bound gives Bayes absolute-error
risk at least
\[
 \int_0^{\ell/(2D)}\left(1-\frac{2Dv}{\ell}\right)dv
 =\frac\ell{4D}.
\]
Averaging the independent seeds preserves it. Consequently some
parameter in the prior support has risk at least
\[
 \frac1{32(\lfloor M/(k+1)\rfloor+1)}
 \ge\frac{k+1}{32(M+k+1)}\ge\frac{k}{96M}.
\]
IID stationarity and $P^2(u,v)=\Pi(v)$ again give the same source
class and supplied certificate, with target $p$ at every row. Combining
the clean floor and edit floor yields $g_M(k)/96$, also at $k=0$.
A predictor with less information can only have a larger minimax risk,
and allowing arbitrary starts contains the stationary experiment.
\end{proof}

\section{Simultaneous finite-horizon competitive factors}
\label{sec:finite-window-proof}

We prove Theorem~\ref{thm:finite-window} using the finite estimates proved
in the preceding two sections. No limit in $M$ or $r$ is taken. Throughout
this section $\ell=1+\log((M+2)/r)\ge1$, the certificate is $1/M$, and
$\Gamma$ denotes the finite-horizon quantities of
\eqref{eq:finite-factor}, not the sequence factors $C$.

\begin{proof}[Proof of Theorem~\ref{thm:finite-window}]
\noindent\textbf{Upper bound, including overlapping windows.}
Set $m=\min\{r-2,\lfloor M/2\rfloor\}$ and
$L=1+\log(M/(m+1))$. Since $r-1\ge r/2$ and
$\lfloor M/2\rfloor+1\ge M/2$,
\[
 \frac{M}{m+1}
 \le\max\left\{\frac{2M}{r},2\right\}
 \le\frac{2(M+2)}r.
\]
The last step uses $r\le M+2$ and covers full or overlapping windows.
Hence $L\le\ell+\log2\le2\ell$. Proposition~\ref{prop:finite-upper}
gives one counts-and-suffix map, independent of $k$, with all-start risk
at most $2^{19}Lg_M(k)\le2^{20}\ell g_M(k)$ for every $k$.
Stationary experiments are a subclass of arbitrary-start experiments,
so taking the infima over the same maps also gives
$\Gamma_{\rm stat}(M,r)\le\Gamma_{\rm all}(M,r)$.

\noindent\textbf{Lower bound and the small-logarithm case.}
Fix any allowable map $A$ and write
$G=\Gamma_{\rm stat}(A;M,r)\le\sqrt M$.
Lemma~\ref{lem:oracle-floor}, at $k=0$, implies $G\ge1/96$ even if
full history were available. If $\ell<64$, then
$1/96\ge2^{-30}\ell$, proving the desired bound. For $\ell\ge64$,
there is nothing to prove when $G\ge2^{-30}\ell$. Suppose instead
\begin{equation}\label{eq:fh-contradiction}
 G<2^{-30}\ell.
\end{equation}
As $r\ge2$ and $M\ge2$, $(M+2)/r\le M$. Thus
$\log M\ge\ell-1\ge63$, and $M>2^{80}$ because $80\log2<56<63$.
Choose the integer budget
\begin{equation}\label{eq:fh-budget}
 k=\left\lceil\max\{128r,M^{3/4}\}\right\rceil.
\end{equation}
This budget depends only on $(M,r)$; the same map $A$ is used at this
budget and at zero.

\noindent\textbf{Every finite lower-domain condition.}
Use the constants from Proposition~\ref{prop:finite-lower}, writing
$A_0=2^{12}$ for its numerical constant to distinguish it from the map:
\[
 c=2^{-16},\qquad t=ck/M,\qquad
 u_0=\lfloor ck/4\rfloor,\qquad
 L_-=\log\frac1{4A_0t}.
\]
The choice of $k$ gives $r\le k/128$, $k\ge M^{3/4}$, and
\[
 \frac M{k^2}\le M^{-1/2}<2^{-40}
   =\frac{A_0c^2}{2^{20}},\qquad
 ck\ge2^{-16}M^{3/4}>2^{44}\ge1024.
\]
Moreover
\[
 \frac rM=(1+2/M)e^{1-\ell}\le2e^{-63}<2^{-79}.
\]
Since $k\le128r+M^{3/4}+1$, this implies
\[
 \frac kM<2^{-72}+2^{-20}+2^{-80}<2^{-19}<\frac18.
\]
In particular $k$ is a legal budget, $k\le M/8$, and
$A_0t=k/(16M)<1/128<1/4$. Together with the original integer
$M\ge64,r\ge2$, these verify every condition in
\eqref{eq:finite-lower-domain}. The windows needed in that proof are
nonoverlapping since $r\le k/128\le M/1024$.
Also $ck\ge1024$ implies $u_0\ge ck/8$, so
\begin{equation}\label{eq:fh-prefactor}
 \rho:=\frac{u_0}{8Mt}\ge\frac1{64}.
\end{equation}

\noindent\textbf{The same map's clean risk.}
Let $S$ be the supremum of the clean risks on the stationary alternatives
in Proposition~\ref{prop:finite-lower}. All qualify at the same supplied
certificate, so $S\le G/\sqrt M$. For $x\ge0$, differentiation shows
that $(1+x)e^{-x/4}$ is maximized at $x=3$, with value
$4e^{-3/4}<4$. Thus $\ell\le1+\log M\le4M^{1/4}$ and
\eqref{eq:fh-contradiction} gives the finite inequalities
\begin{equation}\label{eq:fh-clean}
 S<2^{-28}M^{-1/4}\le2^{-21}M^{-1/4}\le t/32.
\end{equation}
The last step uses $k\ge M^{3/4}$ and $c/32=2^{-21}$.
No asymptotic small-order estimate or distinct clean estimator is used.

\noindent\textbf{A logarithm uniform in both parameters.}
Since $4A_0c=1/4$, $L_-=\log(4M/k)$. Put
$x=\max\{128r,M^{3/4}\}\ge1$; then $k=\lceil x\rceil\le2x$, and
\[
 L_-\ge\min\left\{\log\frac M{64r},\ \log(2M^{1/4})\right\}.
\]
The first term equals
$\ell-1-\log(1+2/M)-\log64$ and is at least
$\ell-2-\log64\ge\ell-8\ge\ell/8$, using
$\log(1+2/M)\le1$, $\log64<6$, and $\ell\ge64$.
The second term is at least $(\ell-1)/4\ge\ell/8$.
Consequently $L_-\ge\ell/8$.

\noindent\textbf{Common-kernel comparison and infimum.}
Proposition~\ref{prop:finite-lower} supplies a qualifying stationary null
and one strict source-independent kernel, preserving the actual final pair,
with
\[
 R_{\rm null}\ge\rho\left[\left(t/16-S\right)L_-\right]_+.
\]
Its full proof above includes source legality, mixture normalization,
window matching, and pointwise budget checks. Here the bracket is positive
by \eqref{eq:fh-clean}. Since $k\ge M^{3/4}\ge\sqrt M$, one has
$g_M(k)=k/M$. Therefore
\[
 G\ge\frac{R_{\rm null}}{g_M(k)}
 \ge\frac{(1/64)(t/32)(\ell/8)}{k/M}
 =2^{-30}\ell,
\]
contradicting \eqref{eq:fh-contradiction}. The lower bound thus holds for
every map and passes to its infimum without requiring attainment.
The witnesses are stationary and hence are also available in the all-start
class. This completes the finite-horizon proof.
\end{proof}

For clarity about the sequence consequences, if $r_M/M$ has positive
liminf, then $\ell_M$ is bounded eventually, and finitely many earlier
horizons have finite factor; the explicit suffix maps give a single
uniform family. If that liminf is zero, $\ell_M$ diverges on a subsequence,
so the finite lower forbids a uniform factor for any family. For fixed
$\alpha$, the finite bounds give the logarithmic order, while the original
fixed-$\alpha$ proof in Proposition~\ref{prop:lower-consequences} retains
its stronger $2^{-26}\log(1/(32\alpha))$ lower constant. Neither argument
claims a logarithmic risk lower at every individual budget.

\clearpage
\bibliographystyle{plainnat}
\bibliography{references}
\end{document}
